% ============================================================================
%  guia_feedback_ie2.tex
%  Guía de respuesta al feedback del jurado (IE1 -> IE2)
%  Proyecto: Heterogeneidad bancaria, seguro interbancario y requerimientos
%            de encaje en el Perú
%  Compilar: pdflatex -> bibtex -> pdflatex x2
% ============================================================================
\documentclass[11pt]{article}

\usepackage[utf8]{inputenc}
\usepackage[T1]{fontenc}
\usepackage[spanish,es-nodecimaldot]{babel}
\usepackage{amsmath,amssymb,amsthm,mathtools}
\usepackage[margin=2.5cm]{geometry}
\usepackage{booktabs,longtable,array}
\usepackage[most]{tcolorbox}
\usepackage[round]{natbib}
\usepackage{enumitem}
\usepackage{hyperref}
\hypersetup{colorlinks=true,linkcolor=blue!40!black,citecolor=blue!40!black,
            urlcolor=blue!40!black}

\bibliographystyle{apalike}

% --- Cajas -----------------------------------------------------------------
\newtcolorbox{jurado}{colback=blue!5,colframe=blue!40!black,
                      title=Comentario del jurado,
                      fonttitle=\bfseries\small, boxrule=0.6pt}
\newtcolorbox{ie2}{colback=black!4,colframe=black!55,
                   title=Compromiso para IE2,
                   fonttitle=\bfseries\small, boxrule=0.6pt}

% --- Macros de estructura --------------------------------------------------
\newcommand{\diag}[1]{\par\noindent\textbf{Diagn\'ostico.}\ #1\par}
\newcommand{\resp}{\par\noindent\textbf{Respuesta.}\par}
\newcommand{\lit}{\par\noindent\textbf{Respaldo en literatura.}\par}

% --- Notación --------------------------------------------------------------
\newcommand{\E}{\mathbb{E}}
\newcommand{\NU}{\mathrm{NU}}
\newcommand{\keff}{\kappa^{\mathrm{eff}}}

\theoremstyle{definition}
\newtheorem{prop}{Proposici\'on}
\newtheorem{obs}{Observaci\'on}

\title{\textbf{Gu\'ia de respuesta al feedback del jurado}\\[2mm]
\large Heterogeneidad bancaria, seguro interbancario y requerimientos de
encaje en el Per\'u\\[1mm]
\normalsize Del Informe de Estudio 1 al Informe de Estudio 2}

\author{Carlos G\'omez Puican \and Fabi\'an M\'endez Su\'arez}
\date{Asesor: Marco Antonio Ortiz Sosa\\[1mm]
Jurado: PhD C\'esar Salinas Depaz \quad PhD Joao Hern\'andez Calder\'on\\[2mm]
Lima, octubre de 2026}

\begin{document}
\maketitle

\begin{abstract}
\noindent Este documento ordena los dieciocho comentarios recibidos en la
sustentaci\'on del Plan de Trabajo y propone, para cada uno, una respuesta
anclada en las ecuaciones del marco anal\'itico. No es un cap\'itulo del IE2:
es el puente entre la cr\'itica del jurado y el contenido que el IE2 debe
contener. Para cada comentario se indica si exige \emph{modelar} estructura
adicional, \emph{justificar} una decisi\'on ya tomada, \emph{calibrar} un
objeto a datos o \emph{posicionar} el trabajo en la literatura; se desarrolla
la respuesta sobre la notaci\'on de la Secci\'on~\ref{sec:notacion}; y se cierra
con un compromiso expl\'icito de qu\'e entra al IE2 y qu\'e queda fuera de
alcance. Toda referencia citada fue cotejada contra su registro editorial; las
que no pudieron verificarse figuran listadas, sin citarse, en el
Ap\'endice~\ref{sec:verificacion}.
\end{abstract}

\tableofcontents
\newpage

% ===========================================================================
\section{C\'omo leer este documento}
% ===========================================================================

El jurado plante\'o dieciocho observaciones. Le\'idas en conjunto, no son
dieciocho pedidos independientes: son seis preguntas de fondo sobre si el
marco es un modelo o una ilustraci\'on. Este documento las agrupa en esos seis
bloques y responde cada comentario con el mismo formato fijo: cita textual,
diagn\'ostico en una l\'inea, respuesta anclada en ecuaciones, respaldo en
literatura verificada y compromiso para el IE2.

Tres advertencias de lectura, que el documento respeta de principio a fin.

\begin{enumerate}[leftmargin=*]
\item \textbf{Utilizaci\'on no es transferencia neta.} El objeto que el marco
mide hoy es $C^u_k/d^u_k$: cobertura esperada recibida por unidad de dep\'osito.
Hablar de subsidio cruzado exige especificar propiedad del fondo, remuneraci\'on,
repago y absorci\'on de p\'erdidas. Esa distinci\'on no se relaja en ninguna
respuesta; donde se promete avanzar hacia transferencias netas (C4) se dice
expl\'icitamente qu\'e estructura adicional hace falta.

\item \textbf{Honestidad sobre lo no resuelto.} Varias respuestas terminan en
``esto no se hace en el IE2''. Eso es deliberado: el valor del plan est\'a en
que sus hip\'otesis son refutables, no en que todas las extensiones est\'en
prometidas.

\item \textbf{Las referencias est\'an verificadas, no recordadas.} Cada
\texttt{\textbackslash cite} de este documento corresponde a una entrada cotejada
contra revista, volumen y p\'aginas. El Ap\'endice~\ref{sec:verificacion}
reporta qu\'e qued\'o sin verificar y por qu\'e; entre ello, una referencia que
la gu\'ia de trabajo atribu\'ia a uno de los miembros del jurado.
\end{enumerate}

% ===========================================================================
\section{Resumen ejecutivo: mapa de los dieciocho comentarios}
% ===========================================================================

\begin{longtable}{@{}p{0.6cm}p{5.6cm}p{1.3cm}p{2.5cm}p{4.2cm}@{}}
\toprule
\textbf{ID} & \textbf{Comentario} & \textbf{Bloque} & \textbf{Tipo} &
\textbf{Estado} \\
\midrule
\endfirsthead
\toprule
\textbf{ID} & \textbf{Comentario} & \textbf{Bloque} & \textbf{Tipo} &
\textbf{Estado} \\
\midrule
\endhead
C1  & Ser muy exigentes con cada ecuaci\'on, aun la m\'as b\'asica & A & JUSTIFICAR & Derivaciones existen; faltaba exponerlas paso a paso (Ap.~\ref{sec:derivaciones}) \\
C13 & Las capas suenan l\'ogicas, como si forzaran las hip\'otesis & A & JUSTIFICAR & Respuesta fuerte: el modelo puede refutarse a s\'i mismo \\
C16 & Decisi\'on \'optima del portafolio ($\omega_k$) & A & MODELAR & $\omega_k$ es el choque, exógeno; la elecci\'on es sobre $u_k$ \\
C17 & Elecci\'on \'optima de $\omega$ con encaje legal distinto, v\'ia utilidad & A & MODELAR & Reformulado como elecci\'on de cartera con $\rho^s\neq\rho^u$ \\
C18 & Elecciones \'optimas de $k$ (teor\'ia) & A/F & JUSTIFICAR & Criterio formal de tipos informativos \\
C5  & Justificar est\'atico vs din\'amico & B & JUSTIFICAR & Tres argumentos; l\'imite reconocido \\
C4  & La decisi\'on de cr\'edito/dep\'ositos es m\'as din\'amica & B & MODELAR & Extensi\'on de dos per\'iodos especificada \\
C2  & Interbancario vs ventanilla dada la overnight (SOFR) & C & MODELAR & Cascada expl\'icita por costo marginal \\
C3  & Tasa de inter\'es real en soles para la decisi\'on de acceso & C & MODELAR/\allowbreak CALIBRAR & El fractil es invariante al deflactor; el TCR entra entre monedas \\
C7  & Entrada / no entrada de bancos al interbancario & C & MODELAR & Margen extensivo con costo fijo de participaci\'on \\
C14 & No hay tasa objetivo; expectativas distintas por tama\~no & C & MODELAR & Oferta end\'ogena con acaparamiento \\
C12 & Dado el pool de recursos, c\'omo modelar lo que falta & C & MODELAR & Mapa completo + cota de max-flow \\
C6  & Imponer costo de financial distress en la minimizaci\'on & D & MODELAR & Rompe $\pi^*_k=c/\kappa$ com\'un entre tipos \\
C11 & Funciones de optimizaci\'on para el encaje voluntario & D & MODELAR & \textbf{Resuelto}; se documenta \\
C8  & Tipo de cambio real y descalce cambiario & E & MODELAR & Hoy ausente por construcci\'on; tres extensiones ordenadas \\
C15 & La capa de monedas introduce muchas diferencias & E & JUSTIFICAR & Seis supuestos, uno por uno, con grados de libertad contados \\
C9  & Justificar categor\'ias y c\'omo definir $K$ & F & JUSTIFICAR/\allowbreak CALIBRAR & Trade-off sesgo--varianza con referencia formal \\
C10 & Elecci\'on de $F$ por tipo; por qu\'e la misma uniforme & F & JUSTIFICAR & \textbf{Resuelto}: misma $G$, distinta escala $\sigma_k$ \\
\bottomrule
\end{longtable}

\noindent\textbf{Prioridad de profundidad.} C13, C14, C8, C6 y C2 reciben el
tratamiento m\'as extenso por ser los de mayor peso cr\'itico. C11 y C10 est\'an
esencialmente resueltos y se documentan sin reabrirse.

% ===========================================================================
\section{Notaci\'on de referencia}\label{sec:notacion}
% ===========================================================================

Toda respuesta se escribe sobre esta notaci\'on, sin alterarla.

\paragraph{Capa 1 --- balance y saldo tras el choque.}
\begin{align}
b + d\,m &= d + E, & m &= \rho + u, \quad u\ge 0, \label{eq:balance}\\
s &= \begin{cases}
d(\omega+u), & \text{encaje no utilizable (NU)}\\
d(\omega+m), & \text{encaje utilizable (U)}
\end{cases} \label{eq:saldo}
\end{align}

\paragraph{Capa 2 --- riesgo idiosincr\'atico.} $e_k\sim F_k$ con media cero;
d\'eficit y excedente esperados
\begin{equation}
S^-_k=\E_e\big[(-s_k)^+\big],\qquad S^+_k=\E_e\big[(s_k)^+\big].
\label{eq:SpSm}
\end{equation}

\paragraph{Capa 3 --- acceso y racionamiento.}
\begin{equation}
\tau_k=1-e^{-\lambda_k},\qquad
\Psi^-_k=\tau_k\min\!\Big(1,\tfrac{P}{Q}\Big),\qquad
P=\sum_k\mu_k\tau_k S^+_k,\quad Q=\sum_k\mu_k\tau_k S^-_k.
\label{eq:interbancario}
\end{equation}

\paragraph{Capa 4 --- monedas y respaldo.}
\begin{align}
\omega^c_k&=\beta_k A^c+e^c_k, &
A^s&=0,\quad
A^u=\begin{cases}-\eta,&\text{prob. }p\\ 0,&\text{prob. }1-p\end{cases}
\label{eq:monedas}\\
\Delta^c(A)&=\sum_k\mu_k\big(1-\Psi^{-,c}_k(A)\big)S^{c-}_k(A), \label{eq:delta}\\
F^u&=\mathbf{1}\{\NU\}\,\nu^u_F\,\rho^u D^u+R^D, &
\varphi(A)&=\min\!\Big(1,\tfrac{F^u}{\Delta^u(A)}\Big). \label{eq:phi}
\end{align}

\paragraph{Incidencia (objeto central).}
\begin{equation}
C^u_k=\E\big[\varphi(A)\,(1-\Psi^{-,u}_k(A))\,S^{u-}_k(A)\big],
\qquad\text{m\'etrica: } \frac{C^u_k}{d^u_k}\ \text{(utilizaci\'on).}
\label{eq:incidencia}
\end{equation}

\paragraph{Colch\'on voluntario end\'ogeno (nota anal\'itica, octubre 2026).}
El banco elige $u_k$ antes del choque resolviendo
\begin{equation}
\min_{u_k\ge 0}\ \mathcal{C}_k(u_k)= c\,u_k+\kappa\,\E\big[(-\varepsilon_k-u_k)^+\big],
\label{eq:newsvendor}
\end{equation}
con condici\'on de primer orden de fractil cr\'itico
\begin{equation}
F_k(-u^*_k)=\frac{c}{\kappa}
\quad\Longleftrightarrow\quad
\pi^*_k\equiv\Pr(s_k<0)=\frac{c}{\kappa},
\label{eq:cpo}
\end{equation}
y cierre de localizaci\'on--escala $\varepsilon_k=\sigma_k\zeta$, $\zeta\sim G$
sim\'etrica y de soporte completo:
\begin{equation}
u^*_k=\sigma_k\,q,\quad q\equiv G^{-1}\!\Big(1-\frac{c}{\kappa}\Big)>0,
\qquad D^*_k=\sigma_k\,\mathcal{D}(q),\quad
\mathcal{D}(q)\equiv\E\big[(\zeta-q)^+\big].
\label{eq:locscale}
\end{equation}
Con acceso heterog\'eneo, el costo marginal \emph{efectivo} del faltante es
\begin{equation}
\keff_k=\psi_k\,\rho^{ib}+(1-\psi_k)\,\kappa,
\qquad
u^*_k=\sigma_k\,G^{-1}\!\Big(1-\frac{c}{\keff_k}\Big).
\label{eq:kappaeff}
\end{equation}
El equilibrio es el punto fijo del mapa $T:t\mapsto t'$ con
$t=\min(1,P/Q)$, recorriendo
$t\to\keff_k\to u^*_k\to (D_k,X_k)\to (P,Q)\to t'$.

\paragraph{Distribuci\'on del choque.} \emph{Baseline} cuantitativo: normal
est\'andar (opcionalmente truncada por abajo), $F_k(x)=G(x/\sigma_k)$, con
$\sigma_k$ calibrada a la volatilidad del crecimiento de dep\'ositos (SBS) y
$G$ com\'un a ambas monedas. La uniforme se conserva solo como caso cerrado de
la capa base.

% ===========================================================================
\section{Bloque A. Rigor y fundamentaci\'on de la decisi\'on \'optima}
% ===========================================================================

\subsection{C1. Exigencia ecuaci\'on por ecuaci\'on}\label{sub:C1}

\begin{jurado}
``Primero, ser recontra \emph{pickies} con cada ecuaci\'on, as\'i sea la m\'as
b\'asica.''
\end{jurado}

\diag{JUSTIFICAR. Las derivaciones existen; lo que falta es exhibirlas con sus
supuestos, sus casos de borde y sus verificaciones de consistencia.}

\resp
La respuesta no es defender las ecuaciones sino someterlas a un protocolo
visible. El Ap\'endice~\ref{sec:derivaciones} deriva, l\'inea por l\'inea, la
restricci\'on de balance \eqref{eq:balance}, el saldo \eqref{eq:saldo} bajo NU y
U, los momentos parciales \eqref{eq:SpSm}, la regla de racionamiento
\eqref{eq:interbancario}, el respaldo \eqref{eq:phi} y la CPO
\eqref{eq:cpo}. Cada derivaci\'on reporta cuatro cosas: de d\'onde sale la
ecuaci\'on, qu\'e supuesto usa, qu\'e ocurre en los bordes y qu\'e se pierde al
escribirla as\'i.

Dos elementos del protocolo merecen anticiparse aqu\'i porque son los que
convierten la exigencia del jurado en una ventaja.

\paragraph{(i) Identidades de consistencia.} El marco admite verificaciones que
deben reportarse expl\'icitamente, porque si fallan el modelo est\'a mal escrito.
Bajo NU y $\E[\varepsilon_k]=0$, por unidad de dep\'osito,
\begin{equation}
X_k - D_k = \E[s_k] = u_k
\qquad\Longrightarrow\qquad X_k = D_k + u_k,
\label{eq:identidad1}
\end{equation}
de modo que excedente y d\'eficit esperados no son objetos libres: su diferencia
es el colch\'on. En el interbancario, la regla proporcional satisface
\begin{equation}
\sum_k \mu_k \Psi^-_k S^-_k
= \min\!\Big(1,\tfrac{P}{Q}\Big)\sum_k\mu_k\tau_k S^-_k
= \min(P,Q),
\label{eq:identidad2}
\end{equation}
es decir, lo efectivamente transferido nunca excede la oferta y agota la
demanda cuando hay recursos suficientes. Y en el respaldo,
\begin{equation}
\sum_k\mu_k C^u_k=\E\big[\varphi(A)\,\Delta^u(A)\big]
=\E\big[\min\big(F^u,\Delta^u(A)\big)\big]\le F^u,
\label{eq:identidad3}
\end{equation}
que es la condici\'on de que el fondo no reparte m\'as de lo que tiene. Estas
tres identidades deben aparecer en el IE2 como \emph{checks} num\'ericos de la
soluci\'on, no solo como \'algebra.

\paragraph{(ii) Convenciones de borde.} Los operadores $(\cdot)^+$ y $\min$
introducen quiebres. El cuadro siguiente fija el tratamiento y, sobre todo,
declara que ninguna convenci\'on altera asignaciones efectivas.

\begin{center}
\begin{tabular}{@{}p{3.2cm}p{3.6cm}p{7.4cm}@{}}
\toprule
\textbf{Caso} & \textbf{Convenci\'on} & \textbf{Justificaci\'on} \\
\midrule
$Q^c(A)=0$ & $\psi^c_k(A)=0$ & No hay demanda que cubrir; el cociente $P/Q$ es
indeterminado pero la asignaci\'on es cero por ambas v\'ias. \\
$\Delta^u(A)=0$ & $\varphi(A)=1$ & No hay faltante residual; $\varphi$ multiplica
a cero en \eqref{eq:incidencia}. \\
$c/\kappa\ge F_k(0)$ & $u^*_k=0$ & Soluci\'on de esquina de \eqref{eq:newsvendor};
con $G$ sim\'etrica, $F_k(0)=1/2$, luego el colch\'on es interior si y solo si
$c<\kappa/2$. \\
$P/Q>1$ & $\Psi^-_k=\tau_k$ & El racionamiento no muerde; el \'unico l\'imite es
el contacto. \\
\bottomrule
\end{tabular}
\end{center}

\lit
No hay literatura que ``apruebe'' una derivaci\'on, pero s\'i la hay que muestra
que las ecuaciones b\'asicas del marco son contabilidad est\'andar de la firma
bancaria y no construcciones \emph{ad hoc}: \citet{baltensperger1980} sistematiza
la elecci\'on conjunta de cartera y colch\'on bajo la restricci\'on de balance, y
\citet{poole1968} plantea la gesti\'on de reservas como un problema de inventario
estoc\'astico con \'optimo de fractil cr\'itico, que es exactamente la estructura
de \eqref{eq:newsvendor}--\eqref{eq:cpo}, sesenta a\~nos antes de este trabajo.

\begin{ie2}
\textbf{Entra:} un ap\'endice de derivaciones paso a paso; las tres identidades
\eqref{eq:identidad1}--\eqref{eq:identidad3} reportadas como verificaciones
num\'ericas de la soluci\'on; el cuadro de convenciones de borde.
\textbf{Queda fuera:} demostraciones de existencia y unicidad en generalidad
m\'axima; la unicidad del punto fijo se verifica num\'ericamente sobre la grilla
de calibraci\'on y se reporta como tal.
\end{ie2}

\subsection{C13. ``Las capas parecen forzar las hip\'otesis''}\label{sub:C13}

\begin{jurado}
``Las capas suenan bastante l\'ogicas, como si estuvieran tratando de forzar las
hip\'otesis.''
\end{jurado}

\diag{JUSTIFICAR. Es el comentario de mayor peso y admite una respuesta
demostrativa, no ret\'orica.}

\resp
La objeci\'on tiene una forma precisa: si cada capa se introduce para que
aparezca el resultado que se quiere, el modelo no prueba nada. Hay tres
respuestas, y la tercera es la que cierra el punto.

\paragraph{(1) Las capas son una secuencia expositiva de un \emph{solo} modelo.}
El objeto de an\'alisis es la capa~4; las capas 1--3 son restricciones de ese
mismo modelo. Formalmente: la capa~1 es el caso $K=1$, $\sigma_k=\sigma$,
$\tau_k=1$, $c\in\{s\}$; la capa~2 relaja $K=1$; la capa~3 relaja $\tau_k=1$;
la capa~4 relaja la moneda \'unica. No hay par\'ametros nuevos entre capas que no
est\'en en la capa~4. Esto es verificable: la lista de par\'ametros del modelo
completo es $\{\rho^c,\ \sigma_k,\ \lambda_k,\ w_k,\ \mu_k,\ \beta_k,\ p,\ \eta,\
\nu^u_F,\ R^D,\ c,\ \rho^{ib},\ r_W,\ \bar\kappa\}$ y cada capa usa un
subconjunto.

\paragraph{(2) Cada dimensi\'on se disciplina con un momento observable, no con
un grado de libertad.} $\rho^c$ se observa (BCRPData); $w_k$ se observa (SBS);
$\sigma_k$ se calibra a la volatilidad del crecimiento de dep\'ositos por clase
(SBS); $c$, $\rho^{ib}$ y $r_W$ se anclan al corredor de tasas
(Secci\'on~\ref{sub:C3}); $\lambda_k$ no es observable y se recorre en grilla, lo
que el PT ya declara. Los par\'ametros genuinamente libres son \emph{tres}:
$\nu^u_F$ y $R^D$ ---la fracci\'on movilizable del fondo y el respaldo adicional en
d\'olares--- y el costo de liquidaci\'on $\bar\kappa$. Los tres se reportan con
bandas de sensibilidad. Un modelo con tres grados de libertad declarados, y con
una dimensi\'on ($\lambda_k$) recorrida en grilla en lugar de ajustada, no es un
modelo que ``acomoda''.

\paragraph{(3) El modelo puede refutar su propia conjetura.} Este es el
argumento decisivo. Con heterogeneidad \emph{solo} en exposici\'on, el cierre de
localizaci\'on--escala \eqref{eq:locscale} da $u^*_k=\sigma_k q$ con $q$ com\'un y
$D^*_k=\sigma_k\mathcal{D}(q)$: el autoseguro reescala pero no invierte el
ordenamiento por $\sigma_k$. El resultado de incidencia ingenuo sobrevive. Pero
basta activar la capa~3 para que deje de sobrevivir. Con acceso heterog\'eneo,
\eqref{eq:kappaeff} implica
\begin{equation}
\lambda_k\downarrow\ \Longrightarrow\ \tau_k\downarrow\ \Longrightarrow\
\psi_k\downarrow\ \Longrightarrow\ \keff_k\uparrow\ \Longrightarrow\
q_k\equiv G^{-1}\!\Big(1-\frac{c}{\keff_k}\Big)\uparrow\ \Longrightarrow\
u^*_k\uparrow,\quad D^*_k=\sigma_k\mathcal{D}(q_k)\downarrow,
\label{eq:cadena}
\end{equation}
puesto que $\mathcal{D}'(q)=-\big(1-G(q)\big)<0$. La utilizaci\'on esperada del
tipo $k$ por unidad de dep\'osito se mueve entonces por dos fuerzas de signo
opuesto:
\begin{equation}
\frac{C^u_k}{d^u_k}\ \propto\
\underbrace{\big(1-\Psi^{-,u}_k\big)}_{\text{peor acceso}\ \Rightarrow\ \uparrow}
\times
\underbrace{\sigma_k\,\mathcal{D}(q_k)}_{\text{peor acceso}\ \Rightarrow\ \downarrow}.
\label{eq:ambiguedad}
\end{equation}
El signo neto es una cuesti\'on cuantitativa, no te\'orica. Es decir: \emph{el
modelo con colch\'on end\'ogeno puede producir el ordenamiento contrario al que la
hip\'otesis H2 sugiere}. Un marco construido para forzar sus hip\'otesis no
exhibe esta propiedad.

Lo mismo vale para H1. El canal asegurador en d\'olares opera a trav\'es de
$\kappa^{u,\mathrm{eff}}>\kappa^{s,\mathrm{eff}}$, que a su vez depende de que
$\varphi$ sea menos el\'astica en d\'olares. Si el respaldo cubriera todos los
estados relevantes aun con $\rho^u=0$, entonces $\varphi\equiv 1$,
$\kappa^{u,\mathrm{eff}}\to\rho^{ib}$ y el motivo asegurador desaparece. H1 tiene
por tanto una condici\'on de refutaci\'on expl\'icita y calibrable: si al calibrar
resulta $\kappa^{u,\mathrm{eff}}\le\kappa^{s,\mathrm{eff}}$, H1 cae.

\begin{obs}
Conviene presentar \eqref{eq:ambiguedad} en la sustentaci\'on como el resultado
principal del avance desde el IE1, no como una complicaci\'on t\'ecnica. Es la
prueba de que las capas disciplinan en lugar de decorar.
\end{obs}

\lit
Tres apoyos. \citet{poole1968} y \citet{baltensperger1980} muestran que la
estructura de decisi\'on usada no se dise\~n\'o para este resultado: es la forma
can\'onica del problema de reservas. \citet{bianchibigio2022} microfundan el
trade-off entre el costo de mantener liquidez y el de conseguirla tras el choque,
que es el par $(c,\kappa)$ de \eqref{eq:newsvendor}. Y
\citet{bonhomme2022} dan el criterio formal por el cual una partici\'on en $K$
tipos es una aproximaci\'on con error acotado, no una elecci\'on libre
(Secci\'on~\ref{sub:C9}).

\begin{ie2}
\textbf{Entra:} la derivaci\'on \eqref{eq:cadena}--\eqref{eq:ambiguedad} como
resultado del cap\'itulo de modelo; la lista expl\'icita de par\'ametros con su
estatus (observado / calibrado / libre); las dos condiciones de refutaci\'on (de
H1 y de H2) enunciadas antes de los resultados num\'ericos.
\textbf{Queda fuera:} un argumento de identificaci\'on causal del encaje
hist\'orico, que el trabajo no reclama.
\end{ie2}

\subsection{C16 y C17. Decisi\'on \'optima de cartera y encaje legal por moneda}
\label{sub:C16}

\begin{jurado}
``Decisi\'on \'optima del portafolio (omega sub k).''\\[1mm]
``Yo hago una elecci\'on \'optima del omega tomando en cuenta que el encaje legal
es diferente y bas\'andose en funciones de utilidad.''
\end{jurado}

\diag{MODELAR, precedido de una desambiguaci\'on de notaci\'on que el IE2 debe
hacer en la primera p\'agina del cap\'itulo de modelo.}

\resp

\paragraph{Desambiguaci\'on.} En \eqref{eq:monedas}, $\omega^c_k$ es el
\emph{choque} de retiros, no una variable de decisi\'on: el banco no elige cu\'anto
le retiran. La notaci\'on invita al malentendido y debe corregirse en el IE2
(propuesta: renombrar el choque como $\xi^c_k$ y reservar $\omega$ para
ponderadores de cartera, o bien decir expl\'icitamente en el texto que $\omega$ es
exógeno). El comentario del jurado es, le\'ido correctamente, un pedido
leg\'itimo: \emph{d\'onde est\'a la decisi\'on \'optima de cartera}.

\paragraph{La decisi\'on de cartera es la decisi\'on de colch\'on.} De
\eqref{eq:balance}, dados $d^c_k$ y $E^c_k$,
\begin{equation}
b^c_k = d^c_k\big(1-\rho^c-u^c_k\big)+E^c_k,
\label{eq:cartera}
\end{equation}
de modo que elegir $u^c_k$ \emph{es} elegir $b^c_k$: la composici\'on de la cartera
entre cr\'edito y liquidez tiene un solo grado de libertad por moneda. No hay una
decisi\'on de portafolio adicional escondida.

\paragraph{De la funci\'on de utilidad a la CPO newsvendor.} El jurado pide que
esa elecci\'on venga de un problema de utilidad. Escr\'ibase el beneficio esperado
del banco del tipo $k$, por moneda, con $r^c_L$ el retorno del cr\'edito, $r^c_R$
la remuneraci\'on del encaje, $r^c_0$ el retorno del colch\'on voluntario y
$r^c_D$ el costo de los dep\'ositos:
\begin{equation}
\Pi^c_k=r^c_L b^c_k+r^c_R\rho^c d^c_k+r^c_0 u^c_k d^c_k-r^c_D d^c_k
-\keff_k\,d^c_k\,\E\big[(-\varepsilon^c_k-u^c_k)^+\big].
\label{eq:utilidad}
\end{equation}
Sustituyendo \eqref{eq:cartera} y derivando respecto de $u^c_k$, los t\'erminos
que no dependen del colch\'on desaparecen y queda, por unidad de dep\'osito,
\begin{equation}
\max_{u^c_k\ge0}\ \Pi^c_k
\quad\Longleftrightarrow\quad
\min_{u^c_k\ge0}\ \underbrace{(r^c_L-r^c_0)}_{\textstyle \equiv\, c^c}u^c_k
+\keff_k\,\E\big[(-\varepsilon^c_k-u^c_k)^+\big],
\label{eq:equivalencia}
\end{equation}
que es exactamente \eqref{eq:newsvendor}. \textbf{Este es el punto central de la
respuesta a C17:} el problema newsvendor no es una forma reducida impuesta, es
la consecuencia de maximizar beneficio sujeto a la restricci\'on de balance, y
$c$ tiene una lectura expl\'icita --- el margen de colocaci\'on neto sacrificado
por mantener liquidez ociosa.

\paragraph{D\'onde entra el encaje legal diferenciado.} El par\'ametro $\rho^c$
entra por dos v\'ias, y conviene separarlas porque tienen implicancias distintas.

\emph{(a) V\'ia de cantidad (siempre activa).} En \eqref{eq:cartera}, un $\rho^c$
mayor desplaza cr\'edito uno a uno. Como $\rho^u\gg\rho^s$ en el Per\'u
(35{,}37\,\% frente a 5{,}56\,\% en julio de 2026), el desplazamiento es un orden
de magnitud mayor en d\'olares. Pero esta v\'ia no toca $u^{*c}_k$: bajo NU,
$\rho^c$ no aparece en \eqref{eq:equivalencia}.

\emph{(b) V\'ia de sustituci\'on (solo bajo U).} Bajo encaje utilizable,
$m^c_k=\rho^c+u^c_k$ y la restricci\'on $u^c_k\ge0$ se convierte en
$m^c_k\ge\rho^c$. El \'optimo es entonces
\begin{equation}
m^{*c}_k=\max\Big\{\rho^c,\ \sigma^c_k\,G^{-1}\!\Big(1-\frac{c^c}{\keff_k}\Big)\Big\},
\label{eq:umbral}
\end{equation}
es decir, el requerimiento legal \emph{muerde} si y solo si excede la liquidez
que el banco habr\'ia elegido por s\'i mismo. Esto genera tres predicciones
contrastables y, de paso, deriva como resultado el ``umbral de encaje'' que el
Anexo~IV del PT listaba como resultado secundario:
\begin{enumerate}[leftmargin=*,itemsep=1pt]
\item Dado que $\rho^u\gg\rho^s$, el requerimiento en d\'olares es mucho m\'as
probablemente vinculante y el de soles m\'as probablemente holgado; la liquidez
voluntaria observada deber\'ia concentrarse en soles.
\item El encaje vinculante sustituye liquidez voluntaria y reduce
mec\'anicamente la probabilidad de faltante por debajo de $c/\keff_k$: la
``prima'' deja de ser elegida por el banco.
\item Los bancos con $\sigma_k$ alto son los menos afectados por el
requerimiento, porque su \'optimo privado ya est\'a por encima; la incidencia del
encaje como restricci\'on recae en los bancos de bajo riesgo de fondeo --- un
resultado distributivo distinto al de la incidencia del seguro, y que conviene
no mezclar con \'el.
\end{enumerate}

\paragraph{Aversi\'on al riesgo (extensi\'on acotada).} Si en lugar de beneficio
esperado se maximiza $\E[\,\mathcal{U}(\Pi^c_k)\,]$ con $\mathcal{U}$ c\'oncava, la
CPO \eqref{eq:cpo} se convierte en una condici\'on de fractil \emph{ponderada por
utilidad marginal}:
\begin{equation}
F_k(-u^*_k)=\frac{c}{\kappa}\cdot
\frac{\E[\mathcal{U}'(\Pi_k)\mid s_k\ge0]}{\E[\mathcal{U}'(\Pi_k)\mid s_k<0]}
\ <\ \frac{c}{\kappa},
\label{eq:aversion}
\end{equation}
porque la utilidad marginal es mayor en los estados de faltante. La aversi\'on al
riesgo eleva el colch\'on \'optimo y, lo que importa para la incidencia, rompe la
igualdad de $\pi^*_k$ entre tipos si los tipos difieren en capitalizaci\'on. Se
propone enunciar \eqref{eq:aversion} y dejar la cuantificaci\'on fuera del IE2.

\lit
\citet{baltensperger1980} y \citet{poole1968} para la equivalencia entre elecci\'on
de cartera y gesti\'on de reservas bajo la restricci\'on de balance;
\citet{bianchibigio2022} para la versi\'on moderna con costos de fondeo ex post;
\citet{froot1993} para la l\'ogica por la cual el objetivo relevante es convexo en
el faltante de fondos internos, que es lo que justifica \eqref{eq:aversion}.

\begin{ie2}
\textbf{Entra:} la desambiguaci\'on notacional de $\omega$; la derivaci\'on
\eqref{eq:utilidad}--\eqref{eq:equivalencia} como secci\'on del cap\'itulo de
modelo; el resultado \eqref{eq:umbral} con sus tres predicciones.
\textbf{Queda fuera:} la calibraci\'on del caso con aversi\'on al riesgo
\eqref{eq:aversion}; se enuncia y se deja como extensi\'on.
\end{ie2}

\subsection{C18. Elecci\'on te\'orica de los tipos $k$}\label{sub:C18}

\begin{jurado}
``Elecciones \'optimas de $k$ (teor\'ia).''
\end{jurado}

\diag{JUSTIFICAR. Hay un criterio formal de qu\'e hace que un tipo sea un tipo, y
el IE2 debe enunciarlo antes de proponer ninguna partici\'on emp\'irica.}

\resp
En el modelo, un banco entra solo a trav\'es de cuatro objetos:
$(\sigma_k,\lambda_k,w_k)$ y su masa $\mu_k$. Esto permite un criterio exacto.

\begin{prop}[Tipos informativos]
Sea $H$ la distribuci\'on conjunta de $(\sigma,\lambda,w)$ en la poblaci\'on de
bancos, ponderada por dep\'ositos. Todos los objetos de equilibrio del modelo
---$P$, $Q$, $t=\min(1,P/Q)$, $\Delta^c(A)$, $\varphi(A)$ y el perfil de
incidencia $\{C^u_k/d^u_k\}$--- son funcionales de $H$ \'unicamente. En
consecuencia, dos bancos son \emph{model-equivalentes} si y solo si comparten
$(\sigma,\lambda,w)$, y una partici\'on en $K$ clases es \emph{informativa} si y
solo si altera el funcional, esto es, si cambia $t$ o el perfil
$\{C^u_k/d^u_k\}$ m\'as all\'a de la banda de sensibilidad reportada.
\end{prop}

\noindent El contenido pr\'actico de la proposici\'on es doble. Primero, variables
que no entran en $(\sigma,\lambda,w)$ ---tama\~no por s\'i mismo, propiedad del
banco, origen del capital--- no definen tipos: definen tipos solo en la medida
en que \emph{predicen} esos tres momentos. Esto responde directamente a la
sospecha de que las categor\'ias se imponen: la categor\'ia no es un supuesto
primitivo, es un resumen de momentos. Segundo, el criterio de parada es
operativo y se puede reportar: se subdivide mientras subdividir mueva el
funcional.

Formalmente, elegir $K$ es discretizar $H$. \citet{bonhomme2022} muestran que el
sesgo de aproximaci\'on de reemplazar una distribuci\'on continua de
heterogeneidad por $K$ puntos de soporte decrece con $K$ a una tasa conocida
bajo suavidad, mientras que el ruido de estimaci\'on de los momentos de cada
clase crece al repartir la muestra; existe por tanto un $K$ interior \'optimo y no
una preferencia por ``m\'as tipos''. Es el argumento que falta en el PT y el que
el jurado especializado en agentes heterog\'eneos va a buscar.

\lit
\citet{bonhomme2022} para la formalizaci\'on del trade-off de discretizaci\'on;
\citet{bonhommemanresa2015} para la posibilidad de \emph{estimar} los grupos del
panel en lugar de imponerlos (v\'ease Secci\'on~\ref{sub:C9});
\citet{acharya2017systemic} para la l\'ogica de que la contribuci\'on relevante de
una instituci\'on se mide por su aporte al d\'eficit esperado del sistema, que es
justamente el funcional de inter\'es.

\begin{ie2}
\textbf{Entra:} la Proposici\'on~1 enunciada en el cap\'itulo de modelo; el criterio
de parada; la justificaci\'on formal de $K$ por referencia a
\citet{bonhomme2022}.
\textbf{Queda fuera:} una demostraci\'on de la tasa de convergencia en este marco
espec\'ifico; se invoca el resultado general.
\end{ie2}

% ===========================================================================
\section{Bloque B. Est\'atico frente a din\'amico}
% ===========================================================================

\subsection{C5. Justificaci\'on del marco est\'atico}\label{sub:C5}

\begin{jurado}
``Justificar el ser est\'atico vs din\'amico.''
\end{jurado}

\diag{JUSTIFICAR, con reconocimiento expl\'icito de lo que el marco est\'atico no
puede contestar.}

\resp
Tres argumentos, en orden de fuerza.

\paragraph{(1) La pregunta es de reparto, no de trayectoria.} El objeto central
es \eqref{eq:incidencia}: qui\'en utiliza la cobertura por unidad de dep\'osito.
Esa es una pregunta de secci\'on transversal condicional a un estado agregado. Un
modelo din\'amico responder\'ia una pregunta distinta ---c\'omo evoluciona la
utilizaci\'on tras un choque--- que no es la que se formula.

\paragraph{(2) La estructura aseguradora es intr\'insecamente de un disparo.} Un
seguro tiene un momento ex ante (se decide la cobertura sin conocer el
siniestro) y uno ex post (se eval\'ua el siniestro). El marco reproduce
exactamente esa secuencia: $u_k$ se elige antes de $\varepsilon_k$ y de $A$;
$\Psi_k$ y $\varphi$ se resuelven despu\'es. A\~nadir per\'iodos no agrega
estructura aseguradora: agrega reconstituci\'on del colch\'on, que es otro
fen\'omeno. La l\'ogica es la misma por la cual el modelo can\'onico de liquidez
bancaria opera con tres fechas y no con horizonte infinito
\citep{diamonddybvig1983,holmstromtirole1998}.

\paragraph{(3) Tractabilidad e identificaci\'on.} El modelo ya combina
heterogeneidad en dos dimensiones, un estado agregado, un mercado OTC con
racionamiento y una restricci\'on de respaldo finito. Un an\'alogo din\'amico
exigir\'ia una distribuci\'on de estados como objeto de equilibrio y har\'ia que
los resultados dependieran de condiciones iniciales no observadas. El costo no
es solo computacional: es de comunicaci\'on de los mecanismos, que es lo que el
trabajo busca aislar.

\paragraph{Lo que el marco est\'atico no puede decir.} Debe decirse sin rodeos:
(i) no puede calcular transferencias netas, porque el aporte al fondo y el uso
de la cobertura ocurren en el mismo instante (ver C4); (ii) no puede hablar de
reconstituci\'on del colch\'on ni de persistencia de la tensi\'on; (iii) no puede
distinguir iliquidez transitoria de insolvencia. Las tres limitaciones ya est\'an
en el alcance declarado del PT y deben repetirse en el IE2.

\lit
\citet{diamonddybvig1983} y \citet{holmstromtirole1998} para la tradici\'on de
modelos de liquidez con pocas fechas;
\citet{allencarlettigale2009} para el caso en que la intervenci\'on del banco
central se analiza en un marco de horizonte corto sin perder el mecanismo.

\begin{ie2}
\textbf{Entra:} una subsecci\'on expl\'icita ``\'ambito temporal del marco'' con los
tres argumentos y las tres limitaciones.
\textbf{Queda fuera:} cualquier afirmaci\'on sobre din\'amica de ajuste.
\end{ie2}

\subsection{C4. La decisi\'on de cr\'edito y dep\'ositos es din\'amica}
\label{sub:C4}

\begin{jurado}
``La decisi\'on de los bancos es m\'as din\'amica con respecto al cr\'edito y los
dep\'ositos, c\'omo explicar eso en el modelo.''
\end{jurado}

\diag{MODELAR. El l\'imite es real y la extensi\'on m\'inima que lo atiende es,
adem\'as, la que habilita el salto de utilizaci\'on a transferencia neta.}

\resp
Hoy la cartera $(b,d,u)$ es fija ex ante y los dep\'ositos son exógenos. La
extensi\'on m\'inima no es un modelo din\'amico completo: es un modelo de
\textbf{dos per\'iodos} que separa el momento del aporte al fondo del momento del
uso de la cobertura.

\paragraph{Estructura propuesta.}
\begin{itemize}[leftmargin=*,itemsep=1pt]
\item $t=0$: el banco elige $(b^c_k,u^c_k)$ sujeto a \eqref{eq:balance} y
constituye el encaje $\rho^c d^c_k$. El costo de oportunidad del aporte es
$(r^c_L-r^c_R)\rho^c d^c_k$.
\item $t=1$: se realizan $A$ y $\varepsilon^c_k$; operan el interbancario
\eqref{eq:interbancario} y el respaldo \eqref{eq:phi}; el banco recibe
cobertura $C^u_k$.
\item $t=2$: se repaga la asistencia recibida a un costo $r_F$ y el cr\'edito
ajusta: $b^c_{k,2}=b^c_{k,1}-\theta\,C^u_k$, con $\theta\ge0$ el grado en que la
asistencia debe devolverse contrayendo activos.
\end{itemize}

\paragraph{Lo que esto permite calcular.} Con esa estructura, la transferencia
neta esperada del tipo $k$ deja de ser indefinible:
\begin{equation}
T^u_k \;=\; \underbrace{\E\big[\varphi(A)(1-\Psi^{-,u}_k)S^{u-}_k\big]}_{\text{cobertura recibida}}
\;-\;\underbrace{(r^u_L-r^u_R)\rho^u d^u_k}_{\text{costo del aporte}}
\;-\;\underbrace{r_F\,\E\big[\varphi(A)(1-\Psi^{-,u}_k)S^{u-}_k\big]}_{\text{repago}}.
\label{eq:transferencia}
\end{equation}
N\'otese que \eqref{eq:transferencia} exige exactamente las tres piezas
contractuales que el PT declara pendientes: propiedad del fondo (qui\'en recibe
$r^u_R$), remuneraci\'on ($r^u_R$) y repago ($r_F$, $\theta$). Es decir, el
modelo de dos per\'iodos no agrega realismo decorativo: agrega el m\'inimo
necesario para que la palabra ``subsidio'' tenga referente. Mientras esas tres
piezas no se fijen, el trabajo debe seguir reportando $C^u_k/d^u_k$ y no $T^u_k$.

\paragraph{Qu\'e no se promete.} Dep\'ositos end\'ogenos exigen una curva de oferta
de dep\'ositos por banco (y, con ella, competencia bancaria), lo que cambia la
naturaleza del trabajo. En IE2 $d^c_k$ permanece exógeno, calibrado a SBS.

\lit
\citet{holmstromtirole1998} para la separaci\'on entre provisi\'on ex ante de
liquidez y su uso ex post, que es exactamente la estructura de dos fechas
propuesta; \citet{bianchibigio2022} para la versi\'on din\'amica completa de
gesti\'on de liquidez bancaria, \'util como contraste de qu\'e se gana y qu\'e se
pierde; \citet{allencarlettigale2009} para el papel de la intervenci\'on del banco
central cuando el mercado interbancario es incompleto.

\begin{ie2}
\textbf{Entra:} la especificaci\'on del modelo de dos per\'iodos y la f\'ormula
\eqref{eq:transferencia} como \emph{agenda declarada}, con la discusi\'on de qu\'e
supuestos contractuales exige.
\textbf{Queda fuera:} la resoluci\'on num\'erica del modelo de dos per\'iodos y los
dep\'ositos end\'ogenos. El IE2 reporta utilizaci\'on, no transferencias netas.
\end{ie2}

% ===========================================================================
\section{Bloque C. Microestructura del mercado interbancario}
% ===========================================================================

\subsection{C2. Interbancario frente a ventanilla dada la tasa overnight}
\label{sub:C2}

\begin{jurado}
``Modelar la din\'amica de decisi\'on del banco entre decidir ir al interbancario
o irse a la ventanilla del BCRP dada la tasa overnight (SOFR).''
\end{jurado}

\diag{MODELAR. La cascada ya existe impl\'icita en $\rho^{ib}<\kappa$; lo que falta
es escribirla como regla de decisi\'on y anclar cada tasa a un observable.}

\resp

\paragraph{La cascada como comparaci\'on de costos marginales.} Ante un faltante
de una unidad, el banco enfrenta tres fuentes, ordenadas por costo:
\begin{equation}
\rho^{ib,c}\ \ <\ \ r^{c}_{W}\ \ <\ \ \bar\kappa^{c},
\label{eq:cascada}
\end{equation}
donde $\rho^{ib,c}$ es el costo de fondearse en el interbancario, $r^c_W$ el de
la ventanilla o respaldo del banco central, y $\bar\kappa^c$ el costo de liquidar
activos. El banco acude al interbancario siempre que logre contraparte, lo que
ocurre con probabilidad $\psi_k=\tau_k\min(1,P/Q)$; en caso contrario recurre a
la ventanilla, y solo si esta no cubre, liquida activos. El costo marginal
\emph{esperado} del faltante es entonces, exactamente,
\begin{equation}
\keff_k=\psi_k\,\rho^{ib,c}+(1-\psi_k)\,\kappa^c,
\qquad
\kappa^c=\underbrace{\varphi^c\,r^c_W+(1-\varphi^c)\,\bar\kappa^c}_{\text{la ventanilla puede no cubrir}},
\label{eq:cascada2}
\end{equation}
que es \eqref{eq:kappaeff} con el segundo escal\'on abierto. Esta es la
contribuci\'on formal que responde a C2: la decisi\'on interbancario--ventanilla
no se modela como una elecci\'on discreta adicional, sino como la esperanza de
una cascada cuyo \'unico elemento estoc\'astico es el acceso.

\paragraph{Y aqu\'i aparece H1, derivada en lugar de supuesta.} En soles,
$\varphi^s\simeq1$ (ventanilla suficientemente el\'astica) y por tanto
$\kappa^s\simeq r^s_W$. En d\'olares, $\varphi^u=\min(1,F^u/\Delta^u)<1$ en los
estados malos, de modo que $\kappa^u=\varphi^u r^u_W+(1-\varphi^u)\bar\kappa^u
>r^u_W$. Luego $\kappa^{u,\mathrm{eff}}>\kappa^{s,\mathrm{eff}}$ a igual
$\sigma$, lo que por \eqref{eq:cpo} reduce la probabilidad de faltante tolerada
en d\'olares y eleva el valor marginal de la reserva requerida como seguro. H1 no
es un supuesto sobre la elasticidad del respaldo: es la consecuencia de
\eqref{eq:cascada2} dada la finitud de $F^u$.

\paragraph{Anclaje de cada tasa a un observable.}
\begin{center}
\begin{tabular}{@{}p{2.2cm}p{5.2cm}p{6.6cm}@{}}
\toprule
\textbf{Objeto} & \textbf{Soles} & \textbf{D\'olares} \\
\midrule
$\rho^{ib,c}$ & Tasa interbancaria overnight MN (BCRPData), dentro del corredor
& Tasa interbancaria overnight ME local; referencia externa: SOFR
\citep{nyfedsofr} m\'as un \emph{spread} de fondeo en d\'olares \\
$r^c_W$ & Tasa de la ventanilla de cr\'edito de regulaci\'on monetaria (techo del
corredor) & Costo de las operaciones de provisi\'on de liquidez en ME del BCRP
(repos de moneda, \emph{swaps}) \\
$c^c=r^c_L-r^c_0$ & Margen de colocaci\'on neto en MN & Margen de colocaci\'on neto
en ME \\
$\bar\kappa^c$ & Costo de liquidaci\'on de activos (banda de sensibilidad) &
\'Idem, m\'as alto por menor profundidad del mercado en ME \\
\bottomrule
\end{tabular}
\end{center}

\paragraph{Una precisi\'on institucional que conviene hacer expl\'icita.} La SOFR
es una tasa \emph{colateralizada} administrada por el Federal Reserve Bank of
New York sobre el mercado de repos de Treasuries \citep{nyfedsofr}. No es el
costo de fondeo en d\'olares de un banco peruano: es el ancla del tramo libre de
riesgo. Usarla sin \emph{spread} subestimar\'ia $\rho^{ib,u}$. M\'as importante
a\'un: en d\'olares el BCRP \emph{no} es emisor, de modo que su ventanilla en ME es
una facilidad con recursos acotados, no un prestamista de \'ultima instancia
el\'astico. Esa asimetr\'ia es el supuesto institucional que sostiene $F^u$ finito
y debe citarse con respaldo (Secci\'on~\ref{sub:C15}).

\lit
\citet{armenterlester2017} para la estructura formal del corredor de tasas y de
c\'omo el acceso desigual a las facilidades del banco central determina d\'onde se
ubica la tasa overnight dentro del corredor; \citet{bechklee2011} para la
evidencia de que, cuando los participantes enfrentan opciones externas
distintas, el mercado overnight se segmenta y el precio efectivo difiere por
participante --- lo que justifica tratar $\rho^{ib}$ como banda y no como punto;
\citet{afonso2015} para la microfundaci\'on de b\'usqueda y negociaci\'on bilateral;
\citet{bianchibigioengel2021} para la demanda precautoria de liquidez en d\'olares
y el rendimiento de conveniencia del d\'olar, que es la contraparte internacional
de $\kappa^{u,\mathrm{eff}}>\kappa^{s,\mathrm{eff}}$.

\begin{ie2}
\textbf{Entra:} la cascada \eqref{eq:cascada}--\eqref{eq:cascada2} con el
segundo escal\'on abierto; la tabla de anclaje de tasas; la derivaci\'on de H1 a
partir de la cascada; la calibraci\'on de $\rho^{ib,c}$ y $r^c_W$ a BCRPData, con
SOFR como referencia externa del tramo en d\'olares.
\textbf{Queda fuera:} un modelo de b\'usqueda y negociaci\'on a lo
\citet{afonso2015}; el trabajo usa una regla reducida de asignaci\'on y lo
declara.
\end{ie2}

\subsection{C3. Tasa de inter\'es real en soles}\label{sub:C3}

\begin{jurado}
``Buscar c\'omo meter la tasa de inter\'es real en soles para la decisi\'on de
acceder al interbancario o no.''
\end{jurado}

\diag{MODELAR/\allowbreak CALIBRAR, con una precisi\'on que conviene no omitir porque es
justamente el tipo de rigor que pide C1.}

\resp

\paragraph{Una observaci\'on exacta antes de la calibraci\'on.} La decisi\'on de
colch\'on depende de la \emph{raz\'on} $c/\keff_k$, no de los niveles. Si todas las
tasas de una misma moneda se deflactan por el mismo factor $(1+\pi^c)$, la
raz\'on es invariante y el fractil \eqref{eq:cpo} no se mueve. En otras palabras:
\emph{dentro} de una moneda, trabajar en t\'erminos nominales o reales no cambia
ning\'un resultado del modelo, porque todos los objetos relevantes son
\emph{spreads} ($c=r_L-r_0$) o razones de costos. Decir esto expl\'icitamente es
la respuesta rigurosa, y evita el error de presentar como sustantivo un cambio
de unidades.

\paragraph{D\'onde s\'i importa el t\'ermino real.} La tasa real s\'i es
indispensable en el \'unico lugar donde se comparan \emph{dos} monedas: el
contraste $\kappa^{u,\mathrm{eff}}$ contra $\kappa^{s,\mathrm{eff}}$, que es el
contenido de H1. All\'i, pasar los costos en d\'olares a unidades comparables
exige la depreciaci\'on esperada:
\begin{equation}
\kappa^{u,\mathrm{eff}}_{\text{(en soles)}}
\;\simeq\; \kappa^{u,\mathrm{eff}}+\E[\hat e],
\qquad
\pi^{*u}_k=\frac{c^u}{\kappa^{u,\mathrm{eff}}},\quad
\pi^{*s}_k=\frac{c^s}{\kappa^{s,\mathrm{eff}}},
\label{eq:real}
\end{equation}
con $\hat e$ la tasa de depreciaci\'on. Es decir: \textbf{C3 y C8 son el mismo
comentario visto desde dos lados}. La comparabilidad entre monedas introduce el
tipo de cambio, y por tanto la pregunta de qu\'e pasa con el descalce
(Secci\'on~\ref{sub:C8}). Presentarlos as\'i, unidos, es m\'as fuerte que responder
cada uno por separado.

\paragraph{Procedimiento concreto para IE2.}
\begin{enumerate}[leftmargin=*,itemsep=1pt]
\item Construir el corredor de tasas en MN a partir de BCRPData
\citep{bcrpdata}: tasa de referencia, tasa de la ventanilla de cr\'edito (techo) y
de los dep\'ositos overnight (piso); tomar $\rho^{ib,s}$ como la interbancaria
overnight observada.
\item Deflactar por expectativas de inflaci\'on a 12 meses (encuesta del BCRP)
para \emph{reportar} el corredor real; documentar que la deflactaci\'on no mueve
$\pi^{*s}_k$.
\item Para la comparaci\'on entre monedas, expresar $\kappa^{u,\mathrm{eff}}$ en
soles seg\'un \eqref{eq:real} y reportar la sensibilidad de H1 al supuesto sobre
$\E[\hat e]$, que es un par\'ametro con incertidumbre genuina.
\end{enumerate}

\lit
\citet{armas2014} para el funcionamiento conjunto del encaje y la tasa de
referencia dentro del esquema de metas de inflaci\'on peruano, que es el anclaje
institucional del corredor; \citet{armenterlester2017} y \citet{bechklee2011}
para la estructura del corredor y la segmentaci\'on.

\begin{ie2}
\textbf{Entra:} la observaci\'on de invarianza al deflactor com\'un (explicitada
como resultado, no omitida); la calibraci\'on del corredor en MN a BCRPData; la
conversi\'on \eqref{eq:real} para la comparaci\'on entre monedas con banda de
sensibilidad sobre $\E[\hat e]$.
\textbf{Queda fuera:} modelar la formaci\'on de expectativas de depreciaci\'on;
$\E[\hat e]$ entra como par\'ametro con rango.
\end{ie2}

\subsection{C7. Entrada y no entrada al mercado interbancario}\label{sub:C7}

\begin{jurado}
``Modelar en el interbancario las posibilidades de entrada como de no entrada de
los bancos.''
\end{jurado}

\diag{MODELAR. Hoy $\tau_k=1-e^{-\lambda_k}$ es contacto probabil\'istico, no
participaci\'on decidida; falta el margen extensivo.}

\resp

\paragraph{Qu\'e es hoy $\tau_k$ y qu\'e no es.} En \eqref{eq:interbancario},
$\tau_k=1-e^{-\lambda_k}$ es la probabilidad de \emph{al menos un contacto} bajo
un proceso de encuentros de intensidad $\lambda_k$ (derivaci\'on en el
Ap\'endice~\ref{sec:derivaciones}). Es un margen intensivo mec\'anico: el banco
siempre quiere entrar, y a veces no encuentra con qui\'en. El comentario del
jurado se\~nala que falta la decisi\'on previa: participar tiene costos
---l\'ineas de cr\'edito, l\'imites de contraparte, colateral, infraestructura
operativa--- que algunos bancos no cubren.

\paragraph{Extensi\'on propuesta.} Sea $\chi_k\ge0$ el costo fijo de participar.
El banco compara dos configuraciones: dentro del mercado, con
$\psi_k=\tau_k t$ y colch\'on $u^{\mathrm{in}}_k$ de \eqref{eq:kappaeff}; fuera,
con $\psi_k=0$, $\keff_k=\kappa$ y colch\'on $u^{\mathrm{out}}_k>u^{\mathrm{in}}_k$.
La ganancia de participar tiene dos componentes:
\begin{equation}
G_k(t)=\underbrace{(\kappa-\rho^{ib})\,\tau_k\,t\,D_k\big(u^{\mathrm{in}}_k\big)}_{\text{cobertura m\'as barata}}
\;+\;\underbrace{c\,\big(u^{\mathrm{out}}_k-u^{\mathrm{in}}_k\big)}_{\text{ahorro de colch\'on}},
\qquad
\text{entra}\iff G_k(t)\ge\chi_k.
\label{eq:entrada}
\end{equation}
El segundo t\'ermino suele olvidarse y es importante: participar no solo abarata
cubrir el faltante, tambi\'en permite cargar menos liquidez ociosa. Ignorarlo
subestima el valor del acceso y, con ello, sobreestima cu\'antos bancos quedan
fuera.

\paragraph{Consecuencia no trivial: complementariedad estrat\'egica.} La decisi\'on
de entrar depende de $t=\min(1,P/Q)$, que a su vez depende de qui\'enes entran.
Cuando el mercado est\'a racionado ($P<Q$), la entrada de un banco excedentario
eleva $P$ y por tanto $t$, aumentando $G_j(t)$ para todos los dem\'as: la
participaci\'on es un complemento estrat\'egico. Esto abre la posibilidad de
equilibrios m\'ultiples ---un equilibrio de mercado delgado y uno de mercado
profundo--- y conecta directamente con C14: un mercado que depende de la
participaci\'on de todos es un mercado que puede apagarse por coordinaci\'on, no
solo por escasez.

\paragraph{Implicancia emp\'irica.} Si $\chi_k$ se correlaciona negativamente con
el tama\~no, el margen extensivo reproduce la estructura n\'ucleo--periferia
documentada en la literatura \citep{craigvonpeter2014} y es coherente con el
hecho estilizado 2 del Anexo~III del PT (correlaci\'on $-0{,}468$ entre tama\~no y
volatilidad). Permite adem\'as un contraste: los bancos fuera del mercado deben
mostrar liquidez contable sistem\'aticamente mayor, que es precisamente
$u^{\mathrm{out}}_k>u^{\mathrm{in}}_k$.

\lit
\citet{craigvonpeter2014} para la estructura n\'ucleo--periferia como resultado de
equilibrio y no como accidente de los datos;
\citet{afonso2014trading} y \citet{afonso2011stressed} para la evidencia de que
las relaciones de contraparte condicionan el acceso, y de que en tensi\'on el
ajuste ocurre sobre todo en \emph{con qui\'en} se negocia;
\citet{afonso2015} para la microfundaci\'on de b\'usqueda;
\citet{cuba2020} para la estructura de interconexiones del sistema peruano.

\begin{ie2}
\textbf{Entra:} la condici\'on de entrada \eqref{eq:entrada}; la discusi\'on de
complementariedad estrat\'egica; el ejercicio de sensibilidad sobre $\chi_k$ en
una grilla, mostrando el efecto sobre $t$ y sobre el perfil de incidencia.
\textbf{Queda fuera:} la selecci\'on entre equilibrios m\'ultiples; se reporta su
existencia y se eval\'ua cada rama por separado.
\end{ie2}

\subsection{C14. Sin tasa objetivo: expectativas heterog\'eneas y acaparamiento}
\label{sub:C14}

\begin{jurado}
``Hay una demanda y una oferta, los fondos son de un tipo; yo tomo la expectativa
diferente que un banco peque\~no, porque no hay una tasa de inter\'es objetivo
interbancaria.''
\end{jurado}

\diag{MODELAR. Es la cr\'itica m\'as profunda al tratamiento actual del
interbancario y, bien respondida, \emph{refuerza} la tesis del trabajo.}

\resp

\paragraph{Qu\'e supone hoy la regla proporcional.} En \eqref{eq:interbancario},
$P=\sum_k\mu_k\tau_k S^+_k$ asume tres cosas: (i) todo banco excedentario presta
\emph{todo} su excedente; (ii) hay un precio \'unico impl\'icito; (iii) nadie
anticipa nada sobre $A$. Las tres son falsas en el sentido que el jurado
se\~nala, y la tercera es la decisiva.

\paragraph{Paso 1: oferta end\'ogena con valor de opci\'on.} Sea $\ell_k\le X_k$ lo
que un banco excedentario efectivamente presta. Retener liquidez tiene valor
porque el propio banco puede necesitarla. Prestar es \'optimo si el precio cubre
el costo esperado de quedarse sin ella:
\begin{equation}
\ell_k>0 \iff \rho^{ib}\ \ge\ \keff_k\cdot\Pr{}_k\big(\text{faltante propio}\mid \mathcal{I}_k\big),
\qquad
P=\sum_k\mu_k\tau_k\,\ell_k\big(\rho^{ib};\mathcal{I}_k\big),
\label{eq:oferta}
\end{equation}
donde $\mathcal{I}_k$ es el conjunto de informaci\'on del banco $k$. La regla
actual es el caso particular $\ell_k=X_k$, que corresponde a
$\Pr_k(\cdot)\to0$.

\paragraph{Paso 2: expectativas heterog\'eneas --- el punto de Joao.} Si el banco
grande intermedia la mayor parte del flujo cambiario, observa una se\~nal
informativa sobre $A^u$ que el banco peque\~no no observa:
$\Pr_{\text{grande}}(A^u=-\eta\mid\mathcal{I})\neq\Pr_{\text{peque\~no}}(\cdot)$.
Por \eqref{eq:oferta}, el banco informado retira oferta \emph{precisamente} en
el estado en que $Q$ es mayor. El resultado es que la tensi\'on
\begin{equation}
t(A)=\min\!\Big(1,\frac{P(A)}{Q(A)}\Big)
\label{eq:tension}
\end{equation}
cae m\'as r\'apido con la severidad de $A$ de lo que el modelo actual predice, y
puede caer a cero de forma discontinua cuando la se\~nal cruza un umbral. El
coaseguro se autodestruye justo cuando se necesita.

\paragraph{Paso 3: la implicancia que conviene subrayar en la sustentaci\'on.}
Este mecanismo \textbf{no debilita la tesis del trabajo: la refuerza}. Si el
coaseguro privado se apaga end\'ogenamente en los estados malos, entonces el
reaseguro p\'ublico ---el respaldo $F^u$--- no es redundante con el interbancario
sino complementario, y su valor marginal es m\'aximo exactamente donde el
mercado falla. Formalmente, $\partial\Delta^u/\partial(\text{severidad})$ se
vuelve m\'as pronunciada, y por \eqref{eq:phi} $\varphi$ cae m\'as, de modo que
$\kappa^{u,\mathrm{eff}}$ sube: H1 se fortalece. El comentario de Joao, llevado
al modelo, produce evidencia a favor de la hip\'otesis que \'el cuestiona.

\paragraph{Paso 4: la indeterminaci\'on del precio, tratada con honestidad.} Sin
tasa objetivo, $\rho^{ib}$ no est\'a fijada por pol\'itica. En los modelos de
b\'usqueda es el resultado de una negociaci\'on bilateral y exhibe dispersi\'on
\citep{afonso2015}; emp\'iricamente se observa dispersi\'on y segmentaci\'on
\citep{afonso2011stressed,bechklee2011,ashcraft2011}. La decisi\'on de modelaci\'on
que se propone para el IE2 es \emph{no} fingir un precio de equilibrio \'unico:
tratar $\rho^{ib}$ como una banda contenida en el corredor de tasas y reportar
los resultados sobre toda la banda. Es menos elegante y mucho m\'as defendible.

\paragraph{Un matiz que conviene traer antes que el jurado.} La evidencia no
respalda un colapso total: \citet{afonso2011stressed} muestran que tras Lehman el
mercado de fondos federales no se congel\'o, sino que se reprici\'o y se reorden\'o
por relaciones. Es decir, el mecanismo correcto no es ``el mercado desaparece''
sino ``el mercado se vuelve caro y selectivo''. En el marco, eso se traduce en
que $t(A)$ cae \emph{y} la asignaci\'on dentro del mercado deja de ser
proporcional: los bancos con relaciones conservan acceso. Reconocer esto evita
sobrevender el mecanismo de acaparamiento.

\lit
\citet{acharyaskeie2011} para el modelo can\'onico de acaparamiento de liquidez por
bancos apalancados con riesgo de refinanciamiento, que en el extremo congela el
mercado; \citet{heider2015} para el mecanismo m\'as cercano al comentario del
jurado: informaci\'on privada sobre el riesgo de contraparte, selecci\'on adversa y
tres reg\'imenes (normal, tensi\'on con spreads altos, quiebre con acaparamiento);
\citet{galeyorulmazer2013} para el resultado de que el equilibrio de mercado es
ineficiente restringido, con demasiada poca liquidez y acaparamiento ineficiente
---que es exactamente el argumento a favor del respaldo p\'ublico---;
\citet{diamondrajan2011} y \citet{brunnermeierpedersen2009} para el canal de temor
a ventas forzadas y espirales de liquidez;
\citet{allencarlettigale2009} para la justificaci\'on de la intervenci\'on del banco
central cuando el interbancario es incompleto;
\citet{bhattacharya1987} para el argumento, anterior y m\'as b\'asico, de que cuando
la liquidez es compartible aparecen incentivos a aprovechar la de los dem\'as, de
modo que el mecanismo colectivo no es gratuito;
\citet{afonso2011stressed} y \citet{ashcraft2011} para el matiz emp\'irico.

\begin{ie2}
\textbf{Entra:} la oferta end\'ogena \eqref{eq:oferta} con un umbral de
acaparamiento dependiente de la se\~nal; la tensi\'on $t(A)$ como objeto
end\'ogeno \eqref{eq:tension}; el tratamiento de $\rho^{ib}$ como banda dentro del
corredor; la discusi\'on expl\'icita de por qu\'e esto refuerza H1.
\textbf{Queda fuera:} un juego de informaci\'on completo con equilibrio bayesiano;
el IE2 implementa la versi\'on de umbral con se\~nal ex\'ogena y reporta
sensibilidad a su precisi\'on.
\end{ie2}

\subsection{C12. El pool de recursos: qu\'e est\'a dentro y qu\'e falta}
\label{sub:C12}

\begin{jurado}
``Dado el pool de recursos del sistema bancario, c\'omo piensan modelar todo eso
que no est\'a a\'un en el modelo.''
\end{jurado}

\diag{MODELAR, en forma de mapa. $P$ y $Q$ \emph{son} el pool; la respuesta \'util
es decir qu\'e captura y en qu\'e direcci\'on se equivoca.}

\resp

\paragraph{Mapa expl\'icito.}
\begin{center}
\begin{tabular}{@{}p{4.3cm}p{4.6cm}p{5.6cm}@{}}
\toprule
\textbf{Pieza} & \textbf{Estatus actual} & \textbf{C\'omo entrar\'ia} \\
\midrule
Excedentes y d\'eficits & Dentro: $P,Q$ en \eqref{eq:interbancario} & --- \\
Racionamiento & Dentro: $\min(1,P/Q)$ proporcional & --- \\
Respaldo finito & Dentro: $\varphi$ en \eqref{eq:phi} & --- \\
Autoseguro end\'ogeno & Dentro: \eqref{eq:kappaeff} y punto fijo & --- \\
Participaci\'on end\'ogena & Fuera & $\ell_k$, $\chi_k$ en \eqref{eq:entrada} (C7) \\
Acaparamiento estrat\'egico & Fuera & $\ell_k(A;\mathcal{I}_k)$ en \eqref{eq:oferta} (C14) \\
Riesgo de contraparte & Fuera & Oferta efectiva $(1-h_{kj})X_k$ con \emph{haircut} por par \\
Colateral & Fuera & Restricci\'on $\ell_k\le \text{colateral disponible}_k$ \\
Red y \emph{tiering} & Fuera (pool \'unico) & Matriz de acceso $a_{kj}$ en lugar de $\tau_k$ \\
Din\'amica & Fuera & Dos per\'iodos (C4) \\
\bottomrule
\end{tabular}
\end{center}

\paragraph{Un resultado que conviene tener a mano.} La simplificaci\'on del pool
no es neutral: se equivoca en una direcci\'on conocida.

\begin{prop}[La regla de pool es una cota superior del coaseguro]
Con acceso heterog\'eneo, la reasignaci\'on factible de liquidez entre bancos es un
problema de flujo m\'aximo sobre el grafo bipartito excedentarios--deficitarios.
La regla proporcional \eqref{eq:interbancario} asigna $\min(P,Q)$, que es el flujo
m\'aximo del grafo \emph{completo}. Como cualquier grafo con menos aristas tiene
flujo m\'aximo no mayor, se sigue que para toda estructura de red
$\sum_k\mu_k\Psi^-_k S^-_k\le\min(P,Q)$. En consecuencia, el marco actual
\textbf{sobreestima} el coaseguro y por tanto \textbf{subestima} el faltante
residual $\Delta^c(A)$ y la utilizaci\'on del respaldo $C^u_k$.
\end{prop}

\noindent La utilidad de esta proposici\'on es doble. Primero, convierte una
limitaci\'on en una cota: los resultados cuantitativos sobre uso del respaldo son
conservadores. Segundo, da un criterio para decidir si vale la pena modelar la
red: si con la cota superior del coaseguro el respaldo ya se agota en los estados
malos, modelarla solo refuerza la conclusi\'on.

\paragraph{Un contrapunto emp\'irico que el jurado puede plantear.}
\citet{cuba2020} encuentran, para el sistema peruano entre 2006 y 2018, que el
riesgo sistémico de liquidez cae en moneda nacional y que no aparece riesgo
significativo en moneda extranjera. Conviene anticiparlo y responderlo: ese
resultado mide contagio realizado de fondeo en la muestra observada, no la
capacidad de respaldo contingente ante un evento agregado $A^u=-\eta$, que por
construcci\'on es poco frecuente. Los dos resultados son compatibles; decirlo
antes de que lo pregunten es mejor que decirlo despu\'es.

\lit
\citet{craigvonpeter2014} para la estructura de red del interbancario;
\citet{rochettirole1996} para la distinci\'on entre el beneficio de coaseguro y el
costo de exposici\'on cruzada; \citet{allen2000contagion} para el resultado de que
las conexiones que comparten riesgo tambi\'en propagan perturbaciones;
\citet{cuba2020} y \citet{carrera2014} para la evidencia y el tratamiento
prudencial del sistema peruano como econom\'ia abierta peque\~na.

\begin{ie2}
\textbf{Entra:} el mapa completo; la Proposici\'on~2 con su implicancia de signo;
la discusi\'on del resultado de \citet{cuba2020}.
\textbf{Queda fuera:} la implementaci\'on de una matriz de acceso $a_{kj}$ y de
riesgo de contraparte; se declaran como agenda y se argumenta por qu\'e la cota
permite postergarlos.
\end{ie2}

% ===========================================================================
\section{Bloque D. Fricciones adicionales en el problema del banco}
% ===========================================================================

\subsection{C6. Costo de financial distress en la minimizaci\'on}\label{sub:C6}

\begin{jurado}
``Al proceso de minimizaci\'on imponerle el costo de \emph{financial distress}.''
\end{jurado}

\diag{MODELAR. No es una adici\'on decorativa: rompe la igualdad
$\pi^*_k=c/\kappa$ entre tipos, que es el supuesto que sostiene la
proporcionalidad del autoseguro.}

\resp

\paragraph{Por qu\'e importa exactamente.} En \eqref{eq:newsvendor} el costo del
faltante es lineal con pendiente $\kappa$. Esa linealidad es lo que produce
\eqref{eq:cpo} con un fractil $c/\kappa$ \emph{com\'un a todos los tipos}, y de
ah\'i $u^*_k\propto\sigma_k$ y $D^*_k\propto\sigma_k$. Si el costo del faltante es
no lineal y su severidad difiere por banco, la proporcionalidad se rompe y el
ordenamiento de incidencia deja de estar determinado por $\sigma_k$. C6 es, en
ese sentido, el segundo mecanismo ---junto con C7 y C14--- capaz de refutar la
conjetura ingenua, y conviene presentarlo as\'i.

\paragraph{Especificaci\'on.} Agr\'eguese un t\'ermino de distress
\begin{equation}
\min_{u_k\ge0}\ \mathcal{C}_k(u_k)=c\,u_k+\kappa\,\E\big[(-\varepsilon_k-u_k)^+\big]
+\Xi_k(u_k),
\qquad
\Xi_k(u_k)=\E\Big[\Lambda_k\big((-\varepsilon_k-u_k)^+\big)\Big],
\label{eq:distress}
\end{equation}
con $\Lambda_k(0)=0$, $\Lambda_k'\ge0$. Dos especificaciones de $\Lambda_k$ dan
resultados cualitativamente distintos y ambas deben discutirse.

\paragraph{Caso (a): distress convexo (ventas forzadas).} Si $\Lambda_k''\ge0$,
el objetivo sigue siendo convexo y la CPO es
\begin{equation}
c=\int_{-\infty}^{-u_k}\big[\kappa+\Lambda_k'(-\varepsilon-u_k)\big]\,dF_k(\varepsilon)
\;>\;\kappa\,F_k(-u_k)
\quad\Longrightarrow\quad
F_k(-u^*_k)<\frac{c}{\kappa}.
\label{eq:cpodistress}
\end{equation}
La probabilidad de faltante tolerada \emph{cae} y el colch\'on \'optimo
\emph{sube}. Si adem\'as $\Lambda_k$ difiere entre tipos ---por valor de
franquicia, por composici\'on de activos, por capitalizaci\'on---, entonces
$\pi^*_k$ ya no es com\'un y $u^*_k$ deja de ser proporcional a $\sigma_k$:
\begin{equation}
u^*_k=\sigma_k\,q_k,\qquad q_k=q\big(c,\kappa,\Lambda_k\big)\ \text{dependiente del tipo.}
\label{eq:qk}
\end{equation}
El cierre de localizaci\'on--escala sobrevive en la forma, pero $q$ deja de ser
com\'un, que era la propiedad que preservaba el ordenamiento.

\paragraph{Caso (b): distress escalonado (corrida o p\'erdida de franquicia).} Si
$\Lambda_k(x)=L_k\mathbf{1}\{x>\theta_k\}$, el objetivo \emph{deja de ser convexo}
y la CPO incorpora un t\'ermino de densidad:
\begin{equation}
c=\kappa F_k(-u_k)+L_k\,f_k(-u_k-\theta_k).
\label{eq:cpoescalon}
\end{equation}
Pueden existir \'optimos locales m\'ultiples y, al variar $L_k$, el \'optimo global
salta de forma discontinua: el banco o bien mantiene un colch\'on peque\~no y
acepta el riesgo de distress, o bien salta a un colch\'on que cubre $\theta_k$. Es
una predicci\'on econ\'omicamente interesante ---no hay autoseguro ``intermedio''
contra la corrida--- y debe reportarse como tal, con la advertencia de que exige
resolver \eqref{eq:distress} por b\'usqueda en grilla y no por CPO.

\paragraph{Interacci\'on con el encaje (y con la lectura macroprudencial).}
Obs\'ervese que $\kappa$ ya era una forma reducida del costo de quedarse corto.
Hacer expl\'icito $\Lambda_k$ separa dos cosas que el PT mantiene unidas: el costo
de \emph{cubrir} el faltante y el costo de \emph{sufrirlo}. Esa separaci\'on le da
al encaje un valor adicional al asegurador: bajo encaje utilizable, elevar
$\rho^c$ sube $m_k$ y reduce $\Pr(\text{distress})$, que es el canal
macroprudencial que el trabajo deliberadamente deja fuera. Conviene nombrarlo y
declararlo fuera de alcance, no ignorarlo.

\lit
\citet{froot1993} para la microfundaci\'on can\'onica: cuando el financiamiento
externo es m\'as caro que el interno, el objetivo es convexo en el faltante de
fondos internos ---es exactamente $\Lambda_k$ convexo---;
\citet{shleifervishny2011} y \citet{diamondrajan2011} para la no linealidad por
ventas forzadas y b\'usqueda de iliquidez;
\citet{brunnermeierpedersen2009} para las espirales entre liquidez de mercado y de
fondeo que generan convexidad;
\citet{acharyaskeie2011} para el canal de apalancamiento y riesgo de
refinanciamiento que hace que $L_k$ var\'ie entre bancos.

\begin{ie2}
\textbf{Entra:} la especificaci\'on \eqref{eq:distress}; el caso convexo con su
CPO \eqref{eq:cpodistress} y el resultado \eqref{eq:qk} de ruptura de la
proporcionalidad; el enunciado del caso escalonado con su advertencia de no
convexidad; la menci\'on expl\'icita del canal macroprudencial como fuera de
alcance.
\textbf{Queda fuera:} la calibraci\'on de $L_k$ y $\theta_k$ a datos peruanos, que
no tiene contraparte observable directa; el caso escalonado se explora solo como
ejercicio de sensibilidad.
\end{ie2}

\subsection{C11. Funciones de optimizaci\'on para el encaje voluntario}
\label{sub:C11}

\begin{jurado}
``Buscar funciones de optimizaci\'on para el encaje voluntario (hecho
parcialmente).''
\end{jurado}

\diag{MODELAR --- \textbf{resuelto}. Corresponde documentar, no reabrir.}

\resp
La nota anal\'itica de octubre de 2026 cierra esta decisi\'on abierta del Anexo~IV
del PT. El contenido a documentar en el IE2 es:

\begin{enumerate}[leftmargin=*,itemsep=2pt]
\item \textbf{El problema.} \eqref{eq:newsvendor}: el banco equilibra el costo de
cargar liquidez ociosa contra el costo esperado de quedarse corto.
\item \textbf{La CPO.} \eqref{eq:cpo}: regla de fractil cr\'itico. Como
$\partial_u\E[(-\varepsilon-u)^+]=-F_k(-u)$, el objetivo es convexo y la
condici\'on es necesaria y suficiente. Lectura econ\'omica: en el \'optimo el banco
tolera una probabilidad de faltante $\pi^*_k=c/\kappa$, id\'entica entre tipos
mientras $\kappa$ sea com\'un.
\item \textbf{El cierre.} \eqref{eq:locscale}: con familia de
localizaci\'on--escala, $u^*_k=\sigma_k q$ y $D^*_k=\sigma_k\mathcal{D}(q)$. Dos
implicancias: el colch\'on es proporcional al riesgo de fondeo ---lo que
racionaliza el hecho estilizado 3 del Anexo~III como \emph{equilibrio} y no como
supuesto--- y la exposici\'on residual tambi\'en lo es, de modo que el autoseguro
reescala pero no invierte el ordenamiento por $\sigma_k$.
\item \textbf{Formas cerradas.} Normal: $u^*_k=\sigma_k z_{1-c/\kappa}$,
$D^*_k=\sigma_k[\phi(q)-q\,c/\kappa]$. Log\'istica:
$u^*_k=b_k\ln(\kappa/c-1)$, $D^*_k=b_k\ln\!\big(\kappa/(\kappa-c)\big)$ con
$b_k=\sqrt3\,\sigma_k/\pi$. Uniforme: $u^*_k=a_k(1-2c/\kappa)$,
$D^*_k=a_k(c/\kappa)^2$ con $a_k=\sqrt3\,\sigma_k$.
\item \textbf{Acceso dentro del autoseguro.} \eqref{eq:kappaeff}: peor acceso
eleva $\keff_k$ y por tanto el colch\'on. Es el canal por el cual la dimensi\'on de
acceso ---y no solo la de exposici\'on--- diferencia a los tipos, y el origen del
argumento de refutabilidad de la Secci\'on~\ref{sub:C13}.
\item \textbf{Equilibrio.} Punto fijo del mapa $T:t\mapsto t'$. Los colchones son
acotados, $T$ es un automapa continuo de $[0,1]$ en s\'i mismo, luego existe punto
fijo por Brouwer. La unicidad no est\'a garantizada y se verifica num\'ericamente
sobre la grilla de calibraci\'on.
\item \textbf{Capa de monedas.} La probabilidad de faltante ex ante en d\'olares
es la mezcla $(1-p)G(-u/\sigma^u_k)+p\,G\big((-u+\beta_k\eta)/\sigma^u_k\big)$,
que no colapsa a un cuantil \'unico; el \'optimo interior es \'unico mientras
$c^u<\kappa^{u,\mathrm{eff}}\Pr(\omega^u_k<0)$ y se resuelve num\'ericamente. El
soporte completo de $G$ es lo que mantiene la mezcla suave; con soporte acotado
la CPO desarrolla tramos planos cuando $\beta_k\eta>2a_k$, raz\'on adicional por
la cual la uniforme es inadecuada aqu\'i.
\end{enumerate}

\noindent \textbf{Lo \'unico pendiente} es la calibraci\'on de $(c,\kappa,\rho^{ib})$
al corredor de tasas (Secci\'on~\ref{sub:C3}), de $R^D$ y de la fracci\'on
movilizable $\nu^u_F$.

\lit
\citet{poole1968} para el antecedente can\'onico de fractil cr\'itico en gesti\'on de
reservas; \citet{baltensperger1980} para la integraci\'on con la elecci\'on de
cartera; \citet{bianchibigio2022} para la microfundaci\'on del par $(c,\kappa)$.

\begin{ie2}
\textbf{Entra:} los siete puntos anteriores como secci\'on del cap\'itulo de modelo,
con la derivaci\'on en el Ap\'endice~\ref{sec:derivaciones}.
\textbf{Queda fuera:} nada de lo resuelto se reabre. La calibraci\'on de
$(c,\kappa,\rho^{ib},R^D,\nu^u_F)$ se reporta con bandas.
\end{ie2}

% ===========================================================================
\section{Bloque E. Dimensi\'on cambiaria}
% ===========================================================================

\subsection{C8. Tipo de cambio real y descalce cambiario}\label{sub:C8}

\begin{jurado}
``El tipo de cambio real y tratar el descalce cambiario.''
\end{jurado}

\diag{MODELAR. La respuesta honesta empieza por reconocer que el descalce no
est\'a impl\'icito en el marco: est\'a \emph{ausente por construcci\'on}.}

\resp

\paragraph{Qu\'e hay hoy, con precisi\'on.} En \eqref{eq:monedas} soles y d\'olares
son circuitos separados: cada moneda tiene su balance
$b^c_k+d^c_k m^c_k=d^c_k+E^c_k$, de modo que \emph{por construcci\'on no hay
descalce}: los activos y pasivos de cada moneda est\'an calzados dentro de su
circuito. Adem\'as no existe tecnolog\'ia de conversi\'on dentro del per\'iodo: el
banco no puede transformar soles en d\'olares, que es justamente lo que hace que
la restricci\'on finita $F^u$ muerda. El tipo de cambio aparece solo de forma
impl\'icita, en la interpretaci\'on del evento $A^u=-\eta$ como salida de
dep\'ositos en ME gatillada por una depreciaci\'on.

Decir esto abiertamente es preferible a sugerir que el TCR est\'a ``incorporado''.
Lo que s\'i puede defenderse es que la ausencia es deliberada: el trabajo aisla el
canal de \emph{cantidad} (liquidez) del canal de \emph{precio} (valuaci\'on), y la
literatura trata el segundo extensamente \citep{cespedes2004,bocolalorenzoni2020}.

\paragraph{Tres extensiones, ordenadas por costo.}

\emph{(E1) Canal de valuaci\'on, versi\'on m\'inima (entra al IE2).} Hacer el evento
agregado expl\'icitamente cambiario: $A^u=-\eta(\hat e)$, con la probabilidad $p$
calibrada a la frecuencia de episodios de depreciaci\'on acompa\~nados de salida de
dep\'ositos en ME, y $\eta$ a su magnitud. No endogeniza el tipo de cambio, pero
le da a $(p,\eta)$ una contraparte emp\'irica en lugar de dejarlos como
par\'ametros libres --- lo que atiende tambi\'en a C13 y C15.

\emph{(E2) El descalce en forma reducida (entra al IE2).} Este es el punto m\'as
\'util y el menos obvio: \textbf{$\beta_k$ ya es, en forma reducida, la exposici\'on
al descalce}. En \eqref{eq:monedas}, $\beta_k$ mide cu\'anto del evento agregado
en d\'olares carga el banco $k$. Un banco con mayor posici\'on neta en ME sufre m\'as
el evento. Luego $\beta_k$ es calibrable con un observable que la SBS publica:
la posici\'on global de cambio por entidad. Esto convierte un par\'ametro que
parec\'ia arbitrario en un momento observado, y responde al jurado con un dato en
lugar de con una promesa.

\emph{(E3) Descalce de balance propiamente dicho (fuera del IE2).} Permitir una
posici\'on neta $x_k=d^u_k-b^u_k\neq0$ y que una depreciaci\'on altere el
patrimonio, $\Delta E_k\simeq -x_k\,\hat e$ expresado en soles. Entonces el choque
de liquidez y el de solvencia se vuelven correlacionados, y el valor asegurador
de $\rho^u$ cambia de signo para bancos largos en d\'olares. Esta es la versi\'on
completa de lo que pide el jurado y conecta el canal de cantidad con el de
precio; tambi\'en cambia la naturaleza del trabajo, porque introduce insolvencia,
que el PT excluye expl\'icitamente de su alcance. Se propone como extensi\'on de
tesis, no de IE2.

\emph{(E4) Mercado de conversi\'on (menci\'on).} Si se permitiera convertir soles en
d\'olares a un costo $\xi$ (spread del \emph{swap} cambiario), $\varphi$ ser\'ia
menos vinculante. Pero $\xi$ sube precisamente en el estado malo: la predicci\'on
del modelo ser\'ia que el valor asegurador de $\rho^u$ aumenta cuando el mercado
de \emph{swaps} se deteriora, que es el fen\'omeno de ``scramble for dollars''
documentado por \citet{bianchibigioengel2021}. Vale la pena mencionarlo como
mecanismo coherente aunque no se implemente.

\lit
\citet{cespedes2004} para el canal de hoja de balance y descalce en econom\'ias
abiertas peque\~nas; \citet{bocolalorenzoni2020} para el prestamista de \'ultima
instancia con recursos limitados en moneda extranjera y para el origen asegurador
de la dolarizaci\'on de pasivos; \citet{dalgic2024} para la lectura del cr\'edito en
ME como arreglo de seguro, que es el contrapunto directo a una lectura puramente
prudencial; \citet{bianchibigioengel2021} para la demanda precautoria de liquidez
en d\'olares; \citet{levyyeyati2008} e \citet{izelevyyeyati2005} para el l\'imite del
prestamista de \'ultima instancia en una econom\'ia dolarizada;
\citet{jeanneranciere2011} para la analog\'ia con las reservas internacionales como
seguro ante interrupciones del financiamiento externo, con la salvedad ---que el
PT ya hace--- de que las reservas soberanas no equivalen al encaje bancario ni
est\'an autom\'aticamente disponibles para cada intermediario.

\begin{ie2}
\textbf{Entra:} el reconocimiento expl\'icito de que no hay descalce en el marco
actual; (E1) la interpretaci\'on cambiaria calibrada de $(p,\eta)$; (E2) la
calibraci\'on de $\beta_k$ a la posici\'on global de cambio por banco (SBS).
\textbf{Queda fuera:} (E3) el descalce de balance con efecto patrimonial y (E4)
el mercado de conversi\'on; se enuncian como agenda de tesis con su mecanismo
esperado.
\end{ie2}

\subsection{C15. Las diferencias que introduce la capa de monedas}
\label{sub:C15}

\begin{jurado}
``Cuando introducen capas, en la de monedas hay mucho que es diferente y debe
sustentarse.''
\end{jurado}

\diag{JUSTIFICAR. La respuesta debe ser exhaustiva: cada diferencia, su
justificaci\'on institucional, su contraparte emp\'irica y qu\'e pasa si se relaja.}

\resp

\begin{longtable}{@{}p{3.5cm}p{4.4cm}p{3.0cm}p{3.6cm}@{}}
\toprule
\textbf{Supuesto} & \textbf{Justificaci\'on} & \textbf{Contraparte emp\'irica} &
\textbf{Si se relaja} \\
\midrule
\endfirsthead
\toprule
\textbf{Supuesto} & \textbf{Justificaci\'on} & \textbf{Contraparte emp\'irica} &
\textbf{Si se relaja} \\
\midrule
\endhead
(i) $A^s=0$ y $A^u=-\eta$ con prob. $p$ &
El BCRP emite soles y no d\'olares: un evento agregado de escasez es posible en ME
y no en MN &
Episodios de salida de dep\'ositos en ME; calibraci\'on de $(p,\eta)$ (C8, E1) &
Un choque agregado en soles ser\'ia absorbido por la ventanilla el\'astica; a\~nade
ruido, no mecanismo \\
(ii) $F^u$ finito, respaldo en soles el\'astico &
No hay prestamista de \'ultima instancia en d\'olares; la capacidad en ME depende de
reservas, no de emisi\'on &
\citet{levyyeyati2008}; \citet{bocolalorenzoni2020} &
\textbf{Es la condici\'on de refutaci\'on de H1}: si $F^u\to\infty$, $\varphi\equiv1$,
$\kappa^{u,\mathrm{eff}}\to\rho^{ib}$ y el motivo asegurador desaparece \\
(iii) $w_k$ fija &
Endogenizarla mezclar\'ia el canal de desdolarizaci\'on con el asegurador &
Dispersi\'on por era, de $\sim$1\,\% a $>$85\,\% (Figura A1 del PT);
\citet{castillo2016} documenta que $w$ responde a pol\'itica &
El encaje diferenciado inducir\'ia sustituci\'on de monedas y el efecto medido
ser\'ia una mezcla de dos canales \\
(iv) $G$ com\'un entre monedas &
La volatilidad idiosincr\'atica de dep\'ositos es aproximadamente sim\'etrica entre
monedas; la asimetr\'ia es agregada y la carga $A^u$ &
\textbf{Contrastable}: comparar $\sigma$ del crecimiento de dep\'ositos MN vs ME
por banco (SBS) &
Si la comparaci\'on rechaza simetr\'ia, se usan $G^s\neq G^u$; los resultados
centrales dependen de simetr\'ia y localizaci\'on--escala, no de normalidad \\
(v) $\rho^c$ uniforme dentro de moneda &
Es el hecho institucional que define la pregunta: el seguro se financia con una
tasa uniforme &
5{,}56\,\% MN y 35{,}37\,\% ME en julio de 2026 (BCRPData);
\citet{cantu2024} halla que los encajes uniformes tienen efecto m\'as fuerte que
los diferenciados &
Con $\rho$ por banco la pregunta distributiva cambia de objeto \\
(vi) $\beta_k$ heterog\'eneo &
No todos los bancos sufren igual el evento agregado; es el descalce en forma
reducida &
Posici\'on global de cambio por entidad (SBS) &
Con $\beta_k$ com\'un, la capa 4 no a\~nade incidencia y solo desplaza el nivel \\
\bottomrule
\end{longtable}

\paragraph{La cuenta de grados de libertad (argumento central).} La capa~4
introduce $\{p,\eta,\beta_k,w_k,\nu^u_F,R^D\}$. De ellos: $w_k$ se observa (SBS);
$\beta_k$ se observa v\'ia posici\'on de cambio (SBS); $(p,\eta)$ se calibran a
episodios documentados. Quedan \textbf{dos} par\'ametros genuinamente libres,
$\nu^u_F$ y $R^D$, y son exactamente los dos sobre los que el trabajo reporta
bandas de sensibilidad. Una capa que a\~nade seis objetos y solo dos grados de
libertad no es una capa que ``acomoda''; presentarlo con esta aritm\'etica es la
forma m\'as econ\'omica de responder simult\'aneamente a C15 y a C13.

\lit
\citet{levyyeyati2008} para el l\'imite del prestamista de \'ultima instancia en una
econom\'ia financieramente dolarizada y la propuesta de requerimientos
descentralizados de liquidez ---el antecedente m\'as directo de leer el encaje en
ME como fondo de seguro---; \citet{bocolalorenzoni2020} para la versi\'on moderna
con recursos limitados en ME; \citet{izelevyyeyati2005} para los determinantes de
la composici\'on monetaria; \citet{castillo2016} para la evidencia peruana de que
$w$ responde a la pol\'itica de encajes; \citet{cantu2024} para el resultado de que
los encajes uniformes tienen efecto m\'as fuerte que los diferenciados;
\citet{dalgic2024} para el contrapunto asegurador. Para la lectura alternativa del
encaje ---la que este trabajo complementa en lugar de disputar---,
\citet{montoromoreno2011} documentan su uso contrac\'iclico en Am\'erica Latina,
\citet{brei2018} lo vinculan a la gesti\'on de flujos de capital y
\citet{glockertowbin2012} condicionan su efectividad a la presencia de fricciones
financieras.

\begin{ie2}
\textbf{Entra:} la tabla completa de supuestos de la capa 4; la cuenta de grados
de libertad; el contraste emp\'irico de simetr\'ia entre monedas del supuesto (iv).
\textbf{Queda fuera:} endogenizar $w_k$, que es un trabajo distinto (y es
exactamente el de \citet{castillo2016}).
\end{ie2}

% ===========================================================================
\section{Bloque F. Heterogeneidad y distribuci\'on}
% ===========================================================================

\subsection{C9. Imponer categor\'ias y definir $K$}\label{sub:C9}

\begin{jurado}
``Justificar la decisi\'on de imponerle una categor\'ia, y c\'omo definir el $K$ de
categor\'ia.''
\end{jurado}

\diag{JUSTIFICAR y CALIBRAR. La Proposici\'on~1 (Secci\'on~\ref{sub:C18}) da el
criterio te\'orico; aqu\'i corresponde el procedimiento emp\'irico.}

\resp

\paragraph{El trade-off, enunciado formalmente.} Elegir $K$ es discretizar la
distribuci\'on $H$ de $(\sigma,\lambda,w)$. Al aumentar $K$: el sesgo de
aproximaci\'on cae ---a una tasa conocida bajo suavidad
\citep{bonhomme2022}---, pero cada clase se estima con menos bancos y el ruido
de sus momentos crece como $1/\sqrt{n_k}$. Existe por tanto un $K$ interior, y la
pregunta ``\,¿por qu\'e no m\'as tipos?\,'' tiene una respuesta t\'ecnica y no una
excusa.

\paragraph{Procedimiento propuesto.}
\begin{enumerate}[leftmargin=*,itemsep=2pt]
\item \textbf{Teor\'ia: $K=2$.} Dos tipos bastan para expresar el mecanismo y
firmar las est\'aticas comparativas. El cap\'itulo de modelo se escribe con
$k\in\{L,S\}$.
\item \textbf{Cuantitativo: $K\in\{2,3,4\}$.} La base de agrupaci\'on son las eras
bancarias de la SBS \citep{sbs} ---segmentaci\'on del propio regulador, lo que evita que la
partici\'on sea una elecci\'on de los autores---, con cortes donde se quiebran las
distribuciones de $(\sigma_k,\text{tama\~no},w_k)$.
\item \textbf{Alternativa disciplinada por datos.} En lugar de imponer los
cortes, estimarlos: los \emph{grouped fixed effects} de
\citet{bonhommemanresa2015} recuperan del panel la partici\'on que mejor explica
la heterogeneidad observada. Reportar ambas y comparar es una respuesta mucho m\'as
fuerte que defender una partici\'on \'unica.
\item \textbf{Criterio de parada.} Subdividir mientras cambie el funcional:
mientras $t=\min(1,P/Q)$ o el perfil $\{C^u_k/d^u_k\}$ se muevan m\'as que la
banda de sensibilidad de $(\nu^u_F,R^D)$. Cuando subdividir no mueva nada, parar.
\item \textbf{Reportar la estabilidad en $K$ como resultado.} Mostrar el perfil
de incidencia para $K=2,3,4$. Si el ordenamiento es estable, eso \emph{es}
evidencia de que el resultado no es un artefacto de la partici\'on; si no lo es,
hay que decirlo.
\end{enumerate}

\paragraph{Lo que no se puede defender.} No se puede afirmar que la
clasificaci\'on por eras SBS identifique tipos econ\'omicos estructurales. Es una
agrupaci\'on administrativa que resulta conveniente porque es exógena a los
autores. El Anexo~III del PT ya advierte que las medias por era se calculan sobre
ventanas desiguales y no prueban persistencia individual; esa advertencia debe
repetirse aqu\'i.

\lit
\citet{bonhomme2022} para el trade-off formal de discretizaci\'on;
\citet{bonhommemanresa2015} para la estimaci\'on de los grupos a partir del panel;
\citet{salinas2024} como referencia metodol\'ogica local de calibraci\'on de un
modelo de agentes heterog\'eneos a momentos de la econom\'ia peruana.

\begin{ie2}
\textbf{Entra:} el procedimiento de cinco pasos; los resultados para $K=2,3,4$;
la comparaci\'on entre partici\'on por eras SBS y partici\'on estimada.
\textbf{Queda fuera:} inferencia sobre el n\'umero de grupos con contrastes
formales; se usa el criterio de estabilidad del funcional.
\end{ie2}

\subsection{C10. Elecci\'on de $F$ por tipo}\label{sub:C10}

\begin{jurado}
``Las elecciones de la funci\'on $F$ por tipo de categor\'ia: por qu\'e asumimos una
uniforme para ambos y no diferentes.''
\end{jurado}

\diag{JUSTIFICAR --- \textbf{resuelto}. Corresponde documentar la decisi\'on
revisada.}

\resp
La premisa del comentario ya no describe al modelo, y conviene decirlo primero:
\textbf{no se asumen distribuciones distintas por tipo, ni se mantiene la
uniforme como \emph{baseline}}. La especificaci\'on revisada es una familia de
localizaci\'on--escala,
\begin{equation}
F_k(x)=G\!\left(\frac{x}{\sigma_k}\right),
\label{eq:locscale2}
\end{equation}
con $G$ est\'andar, sim\'etrica, suave y de soporte completo, \emph{com\'un a todos
los tipos y a ambas monedas}, y $\sigma_k$ calibrada a la volatilidad del
crecimiento de dep\'ositos de cada clase. Es decir, los tipos difieren en escala,
no en forma. Cuatro aclaraciones:

\begin{enumerate}[leftmargin=*,itemsep=2pt]
\item \textbf{Por qu\'e la uniforme sal\'io del \emph{baseline}.} Su soporte acotado
permite que el autoseguro lleve el faltante exactamente a cero y genera quiebres
cuando la capa de monedas desplaza la media: si $\beta_k\eta>2a_k$, la CPO
desarrolla tramos planos. Se conserva solo como caso cerrado de la capa base.
\item \textbf{Por qu\'e normal.} Entrega el mapeo directo de $\sigma_k$ al dise\~no
de mezcla de normales truncadas y la suavidad que hace que $D^c_k(A)$,
$X^c_k(A)$, $\psi^c_k(A)$ y $\varphi(A)$ se comporten bien en los ejercicios de
sensibilidad. Opcionalmente truncada por abajo, porque un drenaje no puede
exceder los dep\'ositos; la truncaci\'on reemplaza $\phi,\Phi$ por momentos
parciales de normal truncada y es inocua para las volatilidades observadas.
\item \textbf{Por qu\'e $G$ com\'un entre monedas.} Porque la volatilidad
idiosincr\'atica de dep\'ositos es aproximadamente sim\'etrica entre monedas; la cola
agregada m\'as gruesa en d\'olares la carga el t\'ermino de evento $A^u$, no el
$\varepsilon$ idiosincr\'atico. Este supuesto es contrastable y el contraste est\'a
comprometido en la Secci\'on~\ref{sub:C15}, supuesto (iv).
\item \textbf{Robustez.} Reemplazar $G$ por una ley sim\'etrica de cola m\'as pesada
($t$ de Student o Laplace) a varianza igualada. Los resultados centrales
---preservaci\'on del ordenamiento por $\sigma_k$ en la exposici\'on residual y el
canal de H1--- dependen solo de simetr\'ia y de la propiedad de
localizaci\'on--escala, y son por tanto independientes de la distribuci\'on.
\end{enumerate}

\lit
\citet{poole1968} y \citet{baltensperger1980} para la tradici\'on en la cual la ley
del choque se mantiene general y la forma cerrada es solo instrumental;
\citet{bonhomme2022} para la l\'ogica paralela de aproximar heterogeneidad con una
estructura disciplinada en lugar de una forma funcional por grupo.

\begin{ie2}
\textbf{Entra:} la especificaci\'on \eqref{eq:locscale2} con $G$ normal como
\emph{baseline}; la uniforme como caso cerrado de la capa base; la robustez a
$t$/Laplace; el contraste de simetr\'ia entre monedas.
\textbf{Queda fuera:} estimar $G$ no param\'etricamente por clase; la muestra de
eras bancarias no lo sostiene.
\end{ie2}

% ===========================================================================
\appendix
\section{Derivaciones paso a paso (respuesta a C1)}\label{sec:derivaciones}
% ===========================================================================

Este ap\'endice deriva, con sus supuestos expl\'icitos, cada ecuaci\'on del marco.

\subsection*{A.1 Restricci\'on de balance}
Los activos del banco son cr\'edito $b$ y reservas $d\,m$; los pasivos son
dep\'ositos $d$ y patrimonio $E$. La identidad contable da
\eqref{eq:balance}: $b+d\,m=d+E$. Despejando,
$b=d(1-m)+E$. \emph{Supuesto:} no hay otras fuentes de fondeo (deuda mayorista,
adeudados), lo que en el caso peruano omite el fondeo externo de algunos bancos
---limitaci\'on que debe declararse---. \emph{Borde:} $m\le1+E/d$, de modo que el
cr\'edito no puede ser negativo.

\subsection*{A.2 Descomposici\'on de las reservas y saldo tras el choque}
Por definici\'on $m=\rho+u$ con $u\ge0$: el encaje exigible m\'as el colch\'on
voluntario. Tras el choque, los dep\'ositos pasan a $d(1+\omega)$, de modo que la
salida de caja es $-d\,\omega$ cuando $\omega<0$. Los recursos l\'iquidos
disponibles para atenderla dependen de la convenci\'on:
\begin{align*}
\text{NU:}&\quad \text{disponible}=d\,u
&\Longrightarrow&\quad s=d\,u+d\,\omega=d(\omega+u),\\
\text{U:}&\quad \text{disponible}=d\,m
&\Longrightarrow&\quad s=d\,m+d\,\omega=d(\omega+m),
\end{align*}
que es \eqref{eq:saldo}. \emph{Interpretaci\'on:} $s<0$ es faltante, $s>0$
excedente. \emph{Supuesto:} el cr\'edito $b$ no es liquidable dentro del per\'iodo
---si lo fuera, su costo de liquidaci\'on es $\bar\kappa$ y entra por
\eqref{eq:cascada}, no por aqu\'i.

\subsection*{A.3 Momentos parciales y su identidad}
Con $D\equiv\E[(-s)^+]$ y $X\equiv\E[(s)^+]$ y la identidad
$(x)^+-(-x)^+=x$,
\[
X-D=\E[s]=d\big(\E[\omega]+u\big)=d\,u
\]
bajo NU y $\E[\omega]=0$, que es \eqref{eq:identidad1}. \emph{Uso:} verificaci\'on
num\'erica obligatoria de la soluci\'on. \emph{Borde:} si $u=0$ y la distribuci\'on
es sim\'etrica, $X=D$.

\subsection*{A.4 Probabilidad de acceso}
Si los encuentros con contrapartes siguen un proceso de Poisson de intensidad
$\lambda_k$ durante el per\'iodo, la probabilidad de \emph{ning\'un} encuentro es
$e^{-\lambda_k}$ y la de al menos uno es $\tau_k=1-e^{-\lambda_k}$.
\emph{Supuesto:} un solo encuentro basta para operar; no hay capacidad limitada
por contraparte. \emph{Borde:} $\lambda_k\to\infty\Rightarrow\tau_k\to1$ (acceso
perfecto); $\lambda_k=0\Rightarrow\tau_k=0$ (banco fuera del mercado, que es el
caso l\'imite de la no participaci\'on de C7).

\subsection*{A.5 Regla de asignaci\'on y su factibilidad}
Con $P=\sum_k\mu_k\tau_k S^+_k$ y $Q=\sum_k\mu_k\tau_k S^-_k$, la regla
$\Psi^-_k=\tau_k\min(1,P/Q)$ transfiere en total
\[
\sum_k\mu_k\Psi^-_k S^-_k=\min\!\Big(1,\frac{P}{Q}\Big)\,Q=\min(P,Q),
\]
que es \eqref{eq:identidad2}: nunca excede la oferta y agota la demanda cuando
$P\ge Q$. \emph{Supuesto:} el racionamiento es proporcional a $\tau_k S^-_k$;
es una construcci\'on propia, motivada por la literatura de gesti\'on de liquidez,
\textbf{no} tomada de \citet{bianchibigio2022} ni de \citet{afonso2015}.
\emph{Borde:} $Q=0\Rightarrow\psi_k=0$ por convenci\'on. \emph{Sesgo conocido:}
por la Proposici\'on~2, la regla es una cota superior del coaseguro factible.

\subsection*{A.6 Respaldo y adding-up}
Con $\Delta^u(A)$ de \eqref{eq:delta} y $\varphi(A)=\min(1,F^u/\Delta^u(A))$,
\[
\sum_k\mu_k C^u_k=\E\big[\varphi(A)\Delta^u(A)\big]
=\E\big[\min(F^u,\Delta^u(A))\big]\le F^u,
\]
que es \eqref{eq:identidad3}. \emph{Supuesto:} el reparto del fondo es
proporcional al faltante residual de cada tipo. \emph{Borde:}
$\Delta^u=0\Rightarrow\varphi=1$ por convenci\'on, sin efecto porque multiplica a
cero. \emph{Nota contable:} bajo encaje utilizable, el principal ya usado por el
banco no vuelve a contarse en el fondo --- de ah\'i el indicador
$\mathbf{1}\{\NU\}$ en \eqref{eq:phi}.

\subsection*{A.7 Convexidad y CPO del colch\'on}
Sea $D_k(u)=\E[(-\varepsilon_k-u)^+]$. Entonces
\[
\frac{\partial D_k}{\partial u}=-F_k(-u),
\qquad
\frac{\partial^2 D_k}{\partial u^2}=f_k(-u)\ \ge 0,
\]
de modo que $\mathcal{C}_k(u)=c\,u+\kappa D_k(u)$ es convexa y la CPO
$c-\kappa F_k(-u)=0$ es necesaria y suficiente, lo que da \eqref{eq:cpo}.
\emph{Borde:} soluci\'on interior si y solo si $c/\kappa<F_k(0)$; con $G$
sim\'etrica, $F_k(0)=1/2$ y la condici\'on es $c<\kappa/2$. En otro caso
$u^*_k=0$. \emph{Est\'aticas comparativas:} de \eqref{eq:locscale},
$\partial u^*_k/\partial\sigma_k=q>0$ y
$\partial u^*_k/\partial(c/\kappa)<0$.

\subsection*{A.8 Cierre de localizaci\'on--escala}
Con $\varepsilon_k=\sigma_k\zeta$ y $\zeta\sim G$ sim\'etrica,
$F_k(-u)=G(-u/\sigma_k)$, de modo que \eqref{eq:cpo} da
$-u^*_k/\sigma_k=G^{-1}(c/\kappa)$, es decir
$u^*_k=\sigma_k G^{-1}(1-c/\kappa)=\sigma_k q$ usando simetr\'ia. Adem\'as
\[
D^*_k=\E\big[(-\sigma_k\zeta-\sigma_k q)^+\big]
=\sigma_k\,\E\big[(-\zeta-q)^+\big]
=\sigma_k\,\E\big[(\zeta-q)^+\big]=\sigma_k\mathcal{D}(q),
\]
donde la \'ultima igualdad usa simetr\'ia de $G$. Y
$\mathcal{D}'(q)=-(1-G(q))<0$, que es la derivada usada en
\eqref{eq:cadena}.

\subsection*{A.9 Existencia del punto fijo}
Sea $t\in[0,1]$ la tensi\'on conjeturada. El mapa recorre
$t\mapsto\psi_k=\tau_k t\mapsto\keff_k$ de \eqref{eq:kappaeff}
$\mapsto u^*_k$ de \eqref{eq:locscale} $\mapsto(D_k,X_k)\mapsto(P,Q)\mapsto
t'=\min(1,P/Q)$. Cada eslab\'on es continuo: $\keff_k$ es af\'in en $\psi_k$;
$u^*_k$ es continuo en $\keff_k$ porque $G^{-1}$ lo es en el interior;
$D_k$ y $X_k$ son continuos en $u_k$; y $\min(1,\cdot)$ es continuo. Los
colchones son acotados, luego $T$ es un automapa continuo de $[0,1]$ en s\'i
mismo y por el teorema de Brouwer tiene punto fijo. \emph{Lo que no se
demuestra:} unicidad. Se verifica num\'ericamente iterando desde m\'ultiples
condiciones iniciales sobre la grilla de calibraci\'on, y se reporta el resultado.

\subsection*{A.10 Capa de monedas: la mezcla}
En d\'olares, $\omega^u_k=\beta_k A^u+\varepsilon^u_k$ con $A^u=-\eta$ con
probabilidad $p$. Luego
\[
\Pr(\omega^u_k<-u)=(1-p)\,G\!\left(\frac{-u}{\sigma^u_k}\right)
+p\,G\!\left(\frac{-u+\beta_k\eta}{\sigma^u_k}\right),
\]
y la CPO es $c^u=\kappa^{u,\mathrm{eff}}\Pr(\omega^u_k<-u^{u*}_k)$. El lado
derecho es estrictamente decreciente en $u$, de modo que el \'optimo interior es
\'unico siempre que $c^u<\kappa^{u,\mathrm{eff}}\Pr(\omega^u_k<0)$; no se reduce a
un cuantil \'unico y se resuelve num\'ericamente. \emph{Por qu\'e importa el soporte
completo:} mantiene las dos componentes de la mezcla solapadas de forma suave;
con soporte acotado la CPO desarrolla tramos planos cuando $\beta_k\eta>2a_k$
---la raz\'on t\'ecnica por la cual la uniforme es inadecuada en esta capa (C10).

% ===========================================================================
\section{Estado de verificaci\'on de las referencias}\label{sec:verificacion}
% ===========================================================================

Toda referencia citada en el cuerpo fue cotejada contra su registro editorial
(revista, volumen, p\'aginas, DOI o serie) en octubre de 2026. El archivo
\texttt{referencias.bib} documenta, entrada por entrada, la fuente del cotejo y
una l\'inea con el hallazgo relevante. Se reportan aqu\'i las excepciones, porque
callarlas ser\'ia el tipo de error que este documento busca evitar.

\paragraph{No verificables, no citadas en el cuerpo.}
\begin{itemize}[leftmargin=*,itemsep=2pt]
\item \textbf{Hern\'andez Calder\'on (2026), \emph{Journal of Economic Dynamics and
Control}.} La b\'usqueda no encontr\'o ning\'un art\'iculo de este autor en esa revista
sobre deuda en moneda extranjera y motivo asegurador. S\'i consta una tesis
doctoral, \emph{Essays in International Macroeconomics and Sovereign Default}
(Washington University in St.\ Louis), cuya defensa figura anunciada para abril
de 2025, con un cap\'itulo titulado ``Partial Default and Exogenous Exchange Rate
Shocks'' listado en SSRN. \textbf{Recomendaci\'on:} pedir la referencia exacta al
propio jurado antes de la sustentaci\'on. Citar el trabajo de un miembro del
jurado con datos equivocados es peor que no citarlo. Para la discusi\'on del
motivo asegurador de la dolarizaci\'on, el documento usa en su lugar
\citet{bocolalorenzoni2020} y \citet{dalgic2024}, que son verificables y cubren
exactamente ese debate.
\end{itemize}

\paragraph{Verificadas con una salvedad.}
\begin{itemize}[leftmargin=*,itemsep=2pt]
\item \textbf{\citet{ortiz2026}.} No verificable en l\'inea. Se conserva porque
figura en la bibliograf\'ia del PT, es la premisa declarada de H1 y su autor es el
asesor del trabajo, de modo que los autores tienen acceso directo al manuscrito.
Confirmar t\'itulo, a\~no y estatus antes de la entrega.
\item \textbf{\citet{heider2015}.} P\'aginas (336--354) y DOI verificados; volumen
y n\'umero no confirmados en el cotejo, por lo que se omiten en la entrada
bibliogr\'afica en lugar de arriesgar un dato falso.
\item \textbf{\citet{castillo2016}.} RePEc y el Reporte de Inflaci\'on del BCRP
listan cuatro autores; la ficha del sitio del BCRP lista tres. Se usa la versi\'on
de cuatro; confirmar con la car\'atula del PDF.
\item \textbf{\citet{bianchibigioengel2021}.} Al cotejo de octubre de 2026 no
consta publicaci\'on en revista; se cita como documento de trabajo del NBER
(n.\,\textordmasculine\ 29457, revisado en marzo de 2023).
\item \textbf{\citet{baltensperger1980}.} La mayor\'ia de cat\'alogos indican
p\'aginas 1--37; una fuente indica 1--38.
\end{itemize}

\paragraph{Atribuci\'on que debe mantenerse.} La regla proporcional
$\tau_k\min(1,P/Q)$ de \eqref{eq:interbancario} es una construcci\'on propia,
motivada por la literatura de gesti\'on de liquidez, y \textbf{no} est\'a tomada de
\citet{bianchibigio2022} ni de \citet{afonso2015}. Esta atribuci\'on debe
mantenerse expl\'icita en el IE2, como ya lo hace el PT.

\bibliography{referencias}

\end{document}
